% Options for packages loaded elsewhere
\PassOptionsToPackage{unicode}{hyperref}
\PassOptionsToPackage{hyphens}{url}
%
\documentclass[
]{article}
\usepackage{amsmath,amssymb}
\usepackage{iftex}
\ifPDFTeX
  \usepackage[T1]{fontenc}
  \usepackage[utf8]{inputenc}
  \usepackage{textcomp} % provide euro and other symbols
\else % if luatex or xetex
  \usepackage{unicode-math} % this also loads fontspec
  \defaultfontfeatures{Scale=MatchLowercase}
  \defaultfontfeatures[\rmfamily]{Ligatures=TeX,Scale=1}
\fi
\usepackage{lmodern}
\ifPDFTeX\else
  % xetex/luatex font selection
\fi
% Use upquote if available, for straight quotes in verbatim environments
\IfFileExists{upquote.sty}{\usepackage{upquote}}{}
\IfFileExists{microtype.sty}{% use microtype if available
  \usepackage[]{microtype}
  \UseMicrotypeSet[protrusion]{basicmath} % disable protrusion for tt fonts
}{}
\makeatletter
\@ifundefined{KOMAClassName}{% if non-KOMA class
  \IfFileExists{parskip.sty}{%
    \usepackage{parskip}
  }{% else
    \setlength{\parindent}{0pt}
    \setlength{\parskip}{6pt plus 2pt minus 1pt}}
}{% if KOMA class
  \KOMAoptions{parskip=half}}
\makeatother
\usepackage{xcolor}
\usepackage[margin=1in]{geometry}
\usepackage{color}
\usepackage{fancyvrb}
\newcommand{\VerbBar}{|}
\newcommand{\VERB}{\Verb[commandchars=\\\{\}]}
\DefineVerbatimEnvironment{Highlighting}{Verbatim}{commandchars=\\\{\}}
% Add ',fontsize=\small' for more characters per line
\usepackage{framed}
\definecolor{shadecolor}{RGB}{248,248,248}
\newenvironment{Shaded}{\begin{snugshade}}{\end{snugshade}}
\newcommand{\AlertTok}[1]{\textcolor[rgb]{0.94,0.16,0.16}{#1}}
\newcommand{\AnnotationTok}[1]{\textcolor[rgb]{0.56,0.35,0.01}{\textbf{\textit{#1}}}}
\newcommand{\AttributeTok}[1]{\textcolor[rgb]{0.13,0.29,0.53}{#1}}
\newcommand{\BaseNTok}[1]{\textcolor[rgb]{0.00,0.00,0.81}{#1}}
\newcommand{\BuiltInTok}[1]{#1}
\newcommand{\CharTok}[1]{\textcolor[rgb]{0.31,0.60,0.02}{#1}}
\newcommand{\CommentTok}[1]{\textcolor[rgb]{0.56,0.35,0.01}{\textit{#1}}}
\newcommand{\CommentVarTok}[1]{\textcolor[rgb]{0.56,0.35,0.01}{\textbf{\textit{#1}}}}
\newcommand{\ConstantTok}[1]{\textcolor[rgb]{0.56,0.35,0.01}{#1}}
\newcommand{\ControlFlowTok}[1]{\textcolor[rgb]{0.13,0.29,0.53}{\textbf{#1}}}
\newcommand{\DataTypeTok}[1]{\textcolor[rgb]{0.13,0.29,0.53}{#1}}
\newcommand{\DecValTok}[1]{\textcolor[rgb]{0.00,0.00,0.81}{#1}}
\newcommand{\DocumentationTok}[1]{\textcolor[rgb]{0.56,0.35,0.01}{\textbf{\textit{#1}}}}
\newcommand{\ErrorTok}[1]{\textcolor[rgb]{0.64,0.00,0.00}{\textbf{#1}}}
\newcommand{\ExtensionTok}[1]{#1}
\newcommand{\FloatTok}[1]{\textcolor[rgb]{0.00,0.00,0.81}{#1}}
\newcommand{\FunctionTok}[1]{\textcolor[rgb]{0.13,0.29,0.53}{\textbf{#1}}}
\newcommand{\ImportTok}[1]{#1}
\newcommand{\InformationTok}[1]{\textcolor[rgb]{0.56,0.35,0.01}{\textbf{\textit{#1}}}}
\newcommand{\KeywordTok}[1]{\textcolor[rgb]{0.13,0.29,0.53}{\textbf{#1}}}
\newcommand{\NormalTok}[1]{#1}
\newcommand{\OperatorTok}[1]{\textcolor[rgb]{0.81,0.36,0.00}{\textbf{#1}}}
\newcommand{\OtherTok}[1]{\textcolor[rgb]{0.56,0.35,0.01}{#1}}
\newcommand{\PreprocessorTok}[1]{\textcolor[rgb]{0.56,0.35,0.01}{\textit{#1}}}
\newcommand{\RegionMarkerTok}[1]{#1}
\newcommand{\SpecialCharTok}[1]{\textcolor[rgb]{0.81,0.36,0.00}{\textbf{#1}}}
\newcommand{\SpecialStringTok}[1]{\textcolor[rgb]{0.31,0.60,0.02}{#1}}
\newcommand{\StringTok}[1]{\textcolor[rgb]{0.31,0.60,0.02}{#1}}
\newcommand{\VariableTok}[1]{\textcolor[rgb]{0.00,0.00,0.00}{#1}}
\newcommand{\VerbatimStringTok}[1]{\textcolor[rgb]{0.31,0.60,0.02}{#1}}
\newcommand{\WarningTok}[1]{\textcolor[rgb]{0.56,0.35,0.01}{\textbf{\textit{#1}}}}
\usepackage{graphicx}
\makeatletter
\newsavebox\pandoc@box
\newcommand*\pandocbounded[1]{% scales image to fit in text height/width
  \sbox\pandoc@box{#1}%
  \Gscale@div\@tempa{\textheight}{\dimexpr\ht\pandoc@box+\dp\pandoc@box\relax}%
  \Gscale@div\@tempb{\linewidth}{\wd\pandoc@box}%
  \ifdim\@tempb\p@<\@tempa\p@\let\@tempa\@tempb\fi% select the smaller of both
  \ifdim\@tempa\p@<\p@\scalebox{\@tempa}{\usebox\pandoc@box}%
  \else\usebox{\pandoc@box}%
  \fi%
}
% Set default figure placement to htbp
\def\fps@figure{htbp}
\makeatother
\setlength{\emergencystretch}{3em} % prevent overfull lines
\providecommand{\tightlist}{%
  \setlength{\itemsep}{0pt}\setlength{\parskip}{0pt}}
\setcounter{secnumdepth}{-\maxdimen} % remove section numbering
\usepackage{bookmark}
\IfFileExists{xurl.sty}{\usepackage{xurl}}{} % add URL line breaks if available
\urlstyle{same}
\hypersetup{
  pdftitle={Financial Econometrics Practice TS \#1},
  hidelinks,
  pdfcreator={LaTeX via pandoc}}

\title{Financial Econometrics Practice TS \#1}
\usepackage{etoolbox}
\makeatletter
\providecommand{\subtitle}[1]{% add subtitle to \maketitle
  \apptocmd{\@title}{\par {\large #1 \par}}{}{}
}
\makeatother
\subtitle{Daniel Campos, Joshua Castillo, \& Cherryl Chico; January 23,
2026}
\author{}
\date{\vspace{-2.5em}}

\begin{document}
\maketitle

\section{1. Gross Domestic Product (GDP) of the United
States}\label{gross-domestic-product-gdp-of-the-united-states}

First of all, let us load the data.

\begin{Shaded}
\begin{Highlighting}[]
\NormalTok{gdpusa }\OtherTok{\textless{}{-}} \FunctionTok{ts}\NormalTok{(}\FunctionTok{read.table}\NormalTok{(}\StringTok{"GDPUSA.dat"}\NormalTok{, }\AttributeTok{header =}\NormalTok{ F), }\AttributeTok{start =} \DecValTok{1947}\NormalTok{, }\AttributeTok{freq =} \DecValTok{4}\NormalTok{)}
\NormalTok{gdpusa}
\end{Highlighting}
\end{Shaded}

\begin{verbatim}
##         Qtr1    Qtr2    Qtr3    Qtr4
## 1947  2033.1  2027.6  2023.5  2055.1
## 1948  2086.0  2120.5  2132.6  2135.0
## 1949  2105.6  2098.4  2120.0  2102.3
## 1950  2184.9  2251.5  2338.5  2383.3
## 1951  2415.7  2457.5  2508.2  2513.7
## 1952  2540.6  2546.0  2564.4  2648.6
## 1953  2697.9  2718.7  2703.4  2662.5
## 1954  2649.8  2652.6  2682.6  2735.1
## 1955  2813.2  2859.0  2897.6  2915.0
## 1956  2903.7  2927.7  2925.0  2973.2
## 1957  2992.2  2985.7  3014.9  2983.7
## 1958  2906.3  2925.4  2993.1  3063.1
## 1959  3121.9  3192.4  3194.7  3203.8
## 1960  3275.8  3258.1  3274.0  3232.0
## 1961  3253.8  3309.1  3372.6  3438.7
## 1962  3500.1  3531.7  3575.1  3586.8
## 1963  3626.0  3666.7  3747.3  3771.8
## 1964  3851.4  3893.3  3954.1  3966.3
## 1965  4062.3  4113.6  4205.1  4302.0
## 1966  4406.7  4421.7  4459.2  4495.8
## 1967  4535.6  4538.4  4581.3  4615.9
## 1968  4710.0  4788.7  4825.8  4844.8
## 1969  4920.6  4935.6  4968.2  4943.9
## 1970  4936.6  4943.6  4989.2  4935.7
## 1971  5069.7  5097.2  5139.1  5151.2
## 1972  5246.0  5365.0  5415.7  5506.4
## 1973  5642.7  5704.1  5674.1  5728.0
## 1974  5678.7  5692.2  5638.4  5616.5
## 1975  5548.2  5587.8  5683.4  5760.0
## 1976  5889.5  5932.7  5965.3  6008.5
## 1977  6079.5  6197.7  6309.5  6309.7
## 1978  6329.8  6574.4  6640.5  6729.8
## 1979  6741.9  6749.1  6799.2  6816.2
## 1980  6837.6  6696.8  6688.8  6813.5
## 1981  6947.0  6895.6  6978.1  6902.1
## 1982  6794.9  6825.9  6799.8  6802.5
## 1983  6892.1  7049.0  7189.9  7339.9
## 1984  7483.4  7612.7  7686.1  7749.2
## 1985  7824.2  7893.1  8013.7  8073.2
## 1986  8148.6  8185.3  8263.6  8308.0
## 1987  8369.9  8460.2  8533.6  8680.2
## 1988  8725.0  8839.6  8891.4  9009.9
## 1989  9101.5  9171.0  9238.9  9257.1
## 1990  9358.3  9392.3  9398.5  9312.9
## 1991  9269.4  9341.6  9388.8  9421.6
## 1992  9534.3  9637.7  9733.0  9834.5
## 1993  9851.0  9908.3  9955.6 10091.0
## 1994 10189.0 10327.0 10387.4 10506.4
## 1995 10543.6 10575.1 10665.1 10737.5
## 1996 10817.9 10998.3 11097.0 11212.2
## 1997 11284.6 11472.1 11615.6 11715.4
## 1998 11832.5 11942.0 12091.6 12287.0
## 1999 12403.3 12498.7 12662.4 12877.6
## 2000 12924.2 13160.8 13178.4 13260.5
## 2001 13222.7 13300.0 13244.8 13280.9
## 2002 13397.0 13478.2 13538.1 13559.0
## 2003 13634.3 13751.5 13985.1 14145.6
## 2004 14221.1 14329.5 14465.0 14609.9
## 2005 14771.6 14839.8 14972.1 15066.6
## 2006 15267.0 15302.7 15326.4 15456.9
## 2007 15493.3 15582.1 15666.7 15762.0
## 2008 15671.4 15752.3 15667.0 15328.0
## 2009 15155.9 15134.1 15189.2 15356.1
## 2010 15415.1 15557.3 15672.0 15750.6
## 2011 15712.8 15825.1 15820.7 16004.1
## 2012 16129.4 16198.8 16220.7 16239.1
## 2013 16383.0 16403.2 16531.7 16663.6
## 2014 16616.5 16841.5 17047.1 17143.0
## 2015 17277.6 17405.7 17463.2 17468.9
## 2016 17556.8 17639.4 17735.1 17824.2
## 2017 17925.3 18021.0 18163.6 18322.5
## 2018 18438.3 18598.1 18732.7 18783.5
## 2019 18927.3 19021.9 19121.1 19219.8
\end{verbatim}

We plot the time series to get a sense of its structure.

\begin{Shaded}
\begin{Highlighting}[]
\FunctionTok{plot}\NormalTok{(gdpusa, }\AttributeTok{xlab =} \StringTok{"Year"}\NormalTok{, }\AttributeTok{ylab =} \StringTok{"US GDP (billions of USD)"}\NormalTok{)}
\end{Highlighting}
\end{Shaded}

\pandocbounded{\includegraphics[keepaspectratio]{CamposCastilloChico_HW1_files/figure-latex/unnamed-chunk-2-1.pdf}}

An overall upward trend in US GDP is clearly visible.

\subsubsection{Is the variance
constant?}\label{is-the-variance-constant}

\emph{Note: Grouping data every two years caused issues with some plots,
so we opted to keep working with the time series as it is.}

We divide the time series in three windows and plot each one in order to
check whether the variance increases over time.

\begin{Shaded}
\begin{Highlighting}[]
\NormalTok{win1 }\OtherTok{\textless{}{-}} \FunctionTok{window}\NormalTok{(gdpusa, }\AttributeTok{start =} \DecValTok{1947}\NormalTok{, }\AttributeTok{end =} \DecValTok{1972}\NormalTok{)}
\NormalTok{win2 }\OtherTok{\textless{}{-}} \FunctionTok{window}\NormalTok{(gdpusa, }\AttributeTok{start =} \DecValTok{1972}\NormalTok{, }\AttributeTok{end =} \DecValTok{1997}\NormalTok{)}
\NormalTok{win3 }\OtherTok{\textless{}{-}} \FunctionTok{window}\NormalTok{(gdpusa, }\AttributeTok{start =} \DecValTok{1997}\NormalTok{, }\AttributeTok{end =} \DecValTok{2019}\NormalTok{)}

\FunctionTok{par}\NormalTok{(}\AttributeTok{mfrow =} \FunctionTok{c}\NormalTok{(}\DecValTok{1}\NormalTok{, }\DecValTok{3}\NormalTok{))}
\FunctionTok{plot}\NormalTok{(win1, }\AttributeTok{ylim =} \FunctionTok{c}\NormalTok{(}\FunctionTok{min}\NormalTok{(gdpusa), }\FunctionTok{max}\NormalTok{(gdpusa)))}
\FunctionTok{plot}\NormalTok{(win2, }\AttributeTok{ylim =} \FunctionTok{c}\NormalTok{(}\FunctionTok{min}\NormalTok{(gdpusa), }\FunctionTok{max}\NormalTok{(gdpusa)))}
\FunctionTok{plot}\NormalTok{(win3, }\AttributeTok{ylim =} \FunctionTok{c}\NormalTok{(}\FunctionTok{min}\NormalTok{(gdpusa), }\FunctionTok{max}\NormalTok{(gdpusa)))}
\end{Highlighting}
\end{Shaded}

\pandocbounded{\includegraphics[keepaspectratio]{CamposCastilloChico_HW1_files/figure-latex/unnamed-chunk-3-1.pdf}}

Although it appears to be the case, we also compute the boxplot.

\begin{Shaded}
\begin{Highlighting}[]
\CommentTok{\# 1. Get the year for every observation in the TS}
\NormalTok{years }\OtherTok{\textless{}{-}} \FunctionTok{floor}\NormalTok{(}\FunctionTok{time}\NormalTok{(gdpusa))}

\CommentTok{\# 2. Create the bins (ensure this is the same length as the data)}
\NormalTok{group\_index }\OtherTok{\textless{}{-}} \DecValTok{2} \SpecialCharTok{*}\NormalTok{ (years }\SpecialCharTok{\%/\%} \DecValTok{2}\NormalTok{)}
\NormalTok{group\_labels }\OtherTok{\textless{}{-}} \FunctionTok{paste0}\NormalTok{(group\_index, }\StringTok{"{-}"}\NormalTok{, group\_index }\SpecialCharTok{+} \DecValTok{1}\NormalTok{)}

\CommentTok{\# 3. Use as.numeric() on the TS to prevent the length mismatch error}
\FunctionTok{boxplot}\NormalTok{(}\FunctionTok{as.numeric}\NormalTok{(gdpusa) }\SpecialCharTok{\textasciitilde{}}\NormalTok{ group\_labels, }
        \AttributeTok{main =} \StringTok{"Quarterly Data by 2{-}Year Intervals"}\NormalTok{,}
        \AttributeTok{xlab =} \StringTok{"Years"}\NormalTok{,}
        \AttributeTok{ylab =} \StringTok{"ln gdpusa"}\NormalTok{)}
\end{Highlighting}
\end{Shaded}

\pandocbounded{\includegraphics[keepaspectratio]{CamposCastilloChico_HW1_files/figure-latex/unnamed-chunk-4-1.pdf}}

Because grouping every two-years still produces a lot of box plots such
that the trend and variance is not yet visible, we will use larger
grouping. We are deciding to group every 10 years. In this graph, it is
more obvious that that the mean is increasing and that the variance is
non-constant, increasing with the level of the time series. A
logarithmic transformation is applied to fix this.

\begin{Shaded}
\begin{Highlighting}[]
\CommentTok{\# 1. Get the year for every observation in the TS}
\NormalTok{years }\OtherTok{\textless{}{-}} \FunctionTok{floor}\NormalTok{(}\FunctionTok{time}\NormalTok{(gdpusa))}

\CommentTok{\# 2. Create the bins (ensure this is the same length as the data)}
\NormalTok{group\_index }\OtherTok{\textless{}{-}} \DecValTok{10} \SpecialCharTok{*}\NormalTok{ (years }\SpecialCharTok{\%/\%} \DecValTok{10}\NormalTok{)}
\NormalTok{group\_labels }\OtherTok{\textless{}{-}} \FunctionTok{paste0}\NormalTok{(group\_index, }\StringTok{"{-}"}\NormalTok{, group\_index }\SpecialCharTok{+} \DecValTok{1}\NormalTok{)}

\CommentTok{\# 3. Use as.numeric() on the TS to prevent the length mismatch error}
\FunctionTok{boxplot}\NormalTok{(}\FunctionTok{as.numeric}\NormalTok{(gdpusa) }\SpecialCharTok{\textasciitilde{}}\NormalTok{ group\_labels, }
        \AttributeTok{main =} \StringTok{"Quarterly Data by 10{-}Year Intervals"}\NormalTok{,}
        \AttributeTok{xlab =} \StringTok{"Years"}\NormalTok{,}
        \AttributeTok{ylab =} \StringTok{"ln gdpusa"}\NormalTok{)}
\end{Highlighting}
\end{Shaded}

\pandocbounded{\includegraphics[keepaspectratio]{CamposCastilloChico_HW1_files/figure-latex/unnamed-chunk-5-1.pdf}}

\begin{Shaded}
\begin{Highlighting}[]
\NormalTok{lngdp }\OtherTok{\textless{}{-}} \FunctionTok{log}\NormalTok{(gdpusa)}
\FunctionTok{plot}\NormalTok{(lngdp)}
\end{Highlighting}
\end{Shaded}

\pandocbounded{\includegraphics[keepaspectratio]{CamposCastilloChico_HW1_files/figure-latex/unnamed-chunk-6-1.pdf}}

The previous plots are drawn again to confirm that the transformation
has indeed stabilized the variance.

\begin{Shaded}
\begin{Highlighting}[]
\NormalTok{lnwin1 }\OtherTok{\textless{}{-}} \FunctionTok{window}\NormalTok{(lngdp, }\AttributeTok{start =} \DecValTok{1947}\NormalTok{, }\AttributeTok{end =} \DecValTok{1972}\NormalTok{)}
\NormalTok{lnwin2 }\OtherTok{\textless{}{-}} \FunctionTok{window}\NormalTok{(lngdp, }\AttributeTok{start =} \DecValTok{1972}\NormalTok{, }\AttributeTok{end =} \DecValTok{1997}\NormalTok{)}
\NormalTok{lnwin3 }\OtherTok{\textless{}{-}} \FunctionTok{window}\NormalTok{(lngdp, }\AttributeTok{start =} \DecValTok{1997}\NormalTok{, }\AttributeTok{end =} \DecValTok{2019}\NormalTok{)}

\FunctionTok{par}\NormalTok{(}\AttributeTok{mfrow =} \FunctionTok{c}\NormalTok{(}\DecValTok{1}\NormalTok{, }\DecValTok{3}\NormalTok{))}
\FunctionTok{plot}\NormalTok{(lnwin1, }\AttributeTok{ylim =} \FunctionTok{c}\NormalTok{(}\FunctionTok{min}\NormalTok{(lngdp), }\FunctionTok{max}\NormalTok{(lngdp)))}
\FunctionTok{plot}\NormalTok{(lnwin2, }\AttributeTok{ylim =} \FunctionTok{c}\NormalTok{(}\FunctionTok{min}\NormalTok{(lngdp), }\FunctionTok{max}\NormalTok{(lngdp)))}
\FunctionTok{plot}\NormalTok{(lnwin3, }\AttributeTok{ylim =} \FunctionTok{c}\NormalTok{(}\FunctionTok{min}\NormalTok{(lngdp), }\FunctionTok{max}\NormalTok{(lngdp)))}
\end{Highlighting}
\end{Shaded}

\pandocbounded{\includegraphics[keepaspectratio]{CamposCastilloChico_HW1_files/figure-latex/unnamed-chunk-7-1.pdf}}

\begin{Shaded}
\begin{Highlighting}[]
\CommentTok{\# 1. Get the year for every observation in the TS}
\NormalTok{years }\OtherTok{\textless{}{-}} \FunctionTok{floor}\NormalTok{(}\FunctionTok{time}\NormalTok{(lngdp))}

\CommentTok{\# 2. Create the bins (ensure this is the same length as the data)}
\NormalTok{group\_index }\OtherTok{\textless{}{-}} \DecValTok{10} \SpecialCharTok{*}\NormalTok{ (years }\SpecialCharTok{\%/\%} \DecValTok{10}\NormalTok{)}
\NormalTok{group\_labels }\OtherTok{\textless{}{-}} \FunctionTok{paste0}\NormalTok{(group\_index, }\StringTok{"{-}"}\NormalTok{, group\_index }\SpecialCharTok{+} \DecValTok{1}\NormalTok{)}

\CommentTok{\# 3. Use as.numeric() on the TS to prevent the length mismatch error}
\FunctionTok{boxplot}\NormalTok{(}\FunctionTok{as.numeric}\NormalTok{(lngdp) }\SpecialCharTok{\textasciitilde{}}\NormalTok{ group\_labels, }
        \AttributeTok{main =} \StringTok{"Quarterly Data by 10{-}Year Intervals"}\NormalTok{,}
        \AttributeTok{xlab =} \StringTok{"Years"}\NormalTok{,}
        \AttributeTok{ylab =} \StringTok{"ln gdpusa"}\NormalTok{)}
\end{Highlighting}
\end{Shaded}

\pandocbounded{\includegraphics[keepaspectratio]{CamposCastilloChico_HW1_files/figure-latex/unnamed-chunk-8-1.pdf}}

The Mean-Variance plot reinforces this. We can disregard the outlier.

\begin{Shaded}
\begin{Highlighting}[]
\NormalTok{mean }\OtherTok{=} \FunctionTok{apply}\NormalTok{(}\FunctionTok{matrix}\NormalTok{(lngdp, }\AttributeTok{ncol =} \DecValTok{10}\NormalTok{), }\DecValTok{2}\NormalTok{, mean)}
\NormalTok{var }\OtherTok{=} \FunctionTok{apply}\NormalTok{(}\FunctionTok{matrix}\NormalTok{(lngdp, }\AttributeTok{ncol =} \DecValTok{10}\NormalTok{), }\DecValTok{2}\NormalTok{, var)}
\FunctionTok{plot}\NormalTok{(var }\SpecialCharTok{\textasciitilde{}}\NormalTok{ mean, }\AttributeTok{main =} \StringTok{"Mean{-}Variance plot"}\NormalTok{)}
\FunctionTok{abline}\NormalTok{(}\FunctionTok{lm}\NormalTok{(var }\SpecialCharTok{\textasciitilde{}}\NormalTok{ mean))}
\end{Highlighting}
\end{Shaded}

\pandocbounded{\includegraphics[keepaspectratio]{CamposCastilloChico_HW1_files/figure-latex/unnamed-chunk-9-1.pdf}}

\subsubsection{Is seasonality present?}\label{is-seasonality-present}

No, as the time series has already been seasonally adjusted.

\subsubsection{Is the mean constant?}\label{is-the-mean-constant}

We can assess whether the mean is constant or not by computing the ACF
and the PACF.

\begin{Shaded}
\begin{Highlighting}[]
\FunctionTok{par}\NormalTok{(}\AttributeTok{mfrow =} \FunctionTok{c}\NormalTok{(}\DecValTok{1}\NormalTok{, }\DecValTok{2}\NormalTok{))}
\FunctionTok{acf}\NormalTok{(lngdp, }\AttributeTok{ylim =} \FunctionTok{c}\NormalTok{(}\SpecialCharTok{{-}}\DecValTok{1}\NormalTok{, }\DecValTok{1}\NormalTok{), }\AttributeTok{lag.max =} \DecValTok{60}\NormalTok{, }\AttributeTok{lwd =} \DecValTok{2}\NormalTok{)}
\FunctionTok{pacf}\NormalTok{(lngdp, }\AttributeTok{ylim =} \FunctionTok{c}\NormalTok{(}\SpecialCharTok{{-}}\DecValTok{1}\NormalTok{, }\DecValTok{1}\NormalTok{), }\AttributeTok{lag.max =} \DecValTok{60}\NormalTok{, }\AttributeTok{lwd =} \DecValTok{2}\NormalTok{)}
\end{Highlighting}
\end{Shaded}

\pandocbounded{\includegraphics[keepaspectratio]{CamposCastilloChico_HW1_files/figure-latex/unnamed-chunk-10-1.pdf}}

As we can observe, the mean is not constant, which goes in line with the
fact that the time series presents an overall upward trend.

Nevertheless, since such trend is linear, it is possible to get rid of
it by differentiating.

\begin{Shaded}
\begin{Highlighting}[]
\FunctionTok{par}\NormalTok{(}\AttributeTok{mfrow =} \FunctionTok{c}\NormalTok{(}\DecValTok{1}\NormalTok{, }\DecValTok{1}\NormalTok{))}
\NormalTok{d1lngdp }\OtherTok{\textless{}{-}} \FunctionTok{diff}\NormalTok{(lngdp)}
\FunctionTok{plot}\NormalTok{(d1lngdp)}
\FunctionTok{abline}\NormalTok{(}\AttributeTok{h =} \DecValTok{0}\NormalTok{)}
\FunctionTok{abline}\NormalTok{(}\AttributeTok{h =} \FunctionTok{mean}\NormalTok{(d1lngdp), }\AttributeTok{col =} \DecValTok{2}\NormalTok{)}
\end{Highlighting}
\end{Shaded}

\pandocbounded{\includegraphics[keepaspectratio]{CamposCastilloChico_HW1_files/figure-latex/unnamed-chunk-11-1.pdf}}

As we can observe, the plot has changed drastically, and now constant
mean (different from 0) is noticeable.

Moreover, comparing variances allows us to analyze if extra differences
are needed.

\begin{Shaded}
\begin{Highlighting}[]
\FunctionTok{formatC}\NormalTok{(}\FunctionTok{var}\NormalTok{(lngdp), }\AttributeTok{format =} \StringTok{"f"}\NormalTok{, }\AttributeTok{digits =} \DecValTok{7}\NormalTok{)}
\end{Highlighting}
\end{Shaded}

\begin{verbatim}
##    V1         
## V1 "0.4445013"
\end{verbatim}

\begin{Shaded}
\begin{Highlighting}[]
\FunctionTok{formatC}\NormalTok{(}\FunctionTok{var}\NormalTok{(d1lngdp), }\AttributeTok{format =} \StringTok{"f"}\NormalTok{, }\AttributeTok{digits =} \DecValTok{7}\NormalTok{)}
\end{Highlighting}
\end{Shaded}

\begin{verbatim}
##    V1         
## V1 "0.0000863"
\end{verbatim}

\begin{Shaded}
\begin{Highlighting}[]
\FunctionTok{formatC}\NormalTok{(}\FunctionTok{var}\NormalTok{(}\FunctionTok{diff}\NormalTok{(d1lngdp)), }\AttributeTok{format =} \StringTok{"f"}\NormalTok{, }\AttributeTok{digits =} \DecValTok{7}\NormalTok{)}
\end{Highlighting}
\end{Shaded}

\begin{verbatim}
##    V1         
## V1 "0.0001105"
\end{verbatim}

The variance decreases substantially after computing the first
difference, but it barely moves from the first to the second. Thus, no
extra differences should be added.

Let us now plot the updated ACF and PACF.

\begin{Shaded}
\begin{Highlighting}[]
\FunctionTok{par}\NormalTok{(}\AttributeTok{mfrow =} \FunctionTok{c}\NormalTok{(}\DecValTok{1}\NormalTok{, }\DecValTok{2}\NormalTok{))}
\FunctionTok{acf}\NormalTok{(d1lngdp, }\AttributeTok{ylim =} \FunctionTok{c}\NormalTok{(}\SpecialCharTok{{-}}\DecValTok{1}\NormalTok{, }\DecValTok{1}\NormalTok{), }\AttributeTok{lag.max =} \DecValTok{60}\NormalTok{, }\AttributeTok{lwd =} \DecValTok{2}\NormalTok{)}
\FunctionTok{pacf}\NormalTok{(d1lngdp, }\AttributeTok{ylim =} \FunctionTok{c}\NormalTok{(}\SpecialCharTok{{-}}\DecValTok{1}\NormalTok{, }\DecValTok{1}\NormalTok{), }\AttributeTok{lag.max =} \DecValTok{60}\NormalTok{, }\AttributeTok{lwd =} \DecValTok{2}\NormalTok{)}
\end{Highlighting}
\end{Shaded}

\pandocbounded{\includegraphics[keepaspectratio]{CamposCastilloChico_HW1_files/figure-latex/unnamed-chunk-13-1.pdf}}

The ACF swiftly decays to 0, meaning that only the very first lags are
significant.

\textbf{The transformed time series is stationary.}

\begin{Shaded}
\begin{Highlighting}[]
\FunctionTok{par}\NormalTok{(}\AttributeTok{mfrow =} \FunctionTok{c}\NormalTok{(}\DecValTok{1}\NormalTok{, }\DecValTok{2}\NormalTok{))}
\FunctionTok{plot}\NormalTok{(gdpusa, }\AttributeTok{main =} \StringTok{"X\_t"}\NormalTok{)}
\FunctionTok{plot}\NormalTok{(d1lngdp, }\AttributeTok{main =} \StringTok{"W\_t"}\NormalTok{)}
\FunctionTok{abline}\NormalTok{(}\AttributeTok{h =} \DecValTok{0}\NormalTok{)}
\end{Highlighting}
\end{Shaded}

\pandocbounded{\includegraphics[keepaspectratio]{CamposCastilloChico_HW1_files/figure-latex/unnamed-chunk-14-1.pdf}}

\subsection{2. Percentage Variation from the Previous Period of the GDP
of the United
States}\label{percentage-variation-from-the-previous-period-of-the-gdp-of-the-united-states}

Firstly, we import the data.

\begin{Shaded}
\begin{Highlighting}[]
\NormalTok{gdpusavar }\OtherTok{\textless{}{-}} \FunctionTok{ts}\NormalTok{(}\FunctionTok{read.table}\NormalTok{(}\StringTok{"gdpusavar.dat"}\NormalTok{), }\AttributeTok{start =} \DecValTok{1990}\NormalTok{, }\AttributeTok{freq =} \DecValTok{4}\NormalTok{)}
\NormalTok{gdpusavar}
\end{Highlighting}
\end{Shaded}

\begin{verbatim}
##      Qtr1 Qtr2 Qtr3 Qtr4
## 1990  4.4  1.5  0.3 -3.6
## 1991 -1.9  3.2  2.0  1.4
## 1992  4.9  4.4  4.0  4.2
## 1993  0.7  2.3  1.9  5.6
## 1994  3.9  5.5  2.4  4.7
## 1995  1.4  1.2  3.5  2.7
## 1996  3.0  6.8  3.6  4.2
## 1997  2.6  6.8  5.1  3.5
## 1998  4.1  3.8  5.1  6.6
## 1999  3.8  3.1  5.3  7.0
## 2000  1.5  7.5  0.5  2.5
## 2001 -1.1  2.4 -1.7  1.1
## 2002  3.5  2.4  1.8  0.6
## 2003  2.2  3.5  7.0  4.7
## 2004  2.2  3.1  3.8  4.1
## 2005  4.5  1.9  3.6  2.6
## 2006  5.4  0.9  0.6  3.5
## 2007  0.9  2.3  2.2  2.5
## 2008 -2.3  2.1 -2.1 -8.4
## 2009 -4.4 -0.6  1.5  4.5
## 2010  1.5  3.7  3.0  2.0
## 2011 -1.0  2.9 -0.1  4.7
## 2012  3.2  1.7  0.5  0.5
## 2013  3.6  0.5  3.2  3.2
## 2014 -1.1  5.5  5.0  2.3
## 2015  3.2  3.0  1.3  0.1
## 2016  2.0  1.9  2.2  2.0
## 2017  2.3  2.2  3.2  3.5
## 2018  2.5  3.5  2.9  1.1
## 2019  3.1  2.0  2.1  2.1
\end{verbatim}

We plot the time series so as to visualize its shape.

\begin{Shaded}
\begin{Highlighting}[]
\FunctionTok{plot}\NormalTok{(gdpusavar, }\AttributeTok{xlab =} \StringTok{"Year"}\NormalTok{, }\AttributeTok{ylab =} \StringTok{"US GDP (quarterly change in \%)"}\NormalTok{, }\AttributeTok{type =} \StringTok{"o"}\NormalTok{)}
\FunctionTok{abline}\NormalTok{(}\AttributeTok{h =} \FunctionTok{mean}\NormalTok{(gdpusavar), }\AttributeTok{col =} \DecValTok{2}\NormalTok{)}
\end{Highlighting}
\end{Shaded}

\pandocbounded{\includegraphics[keepaspectratio]{CamposCastilloChico_HW1_files/figure-latex/unnamed-chunk-16-1.pdf}}

In this case, since data is on the quarterly rate of change of US GDP,
we observe a time series that fluctuates between positive and negative
values, with the biggest outliers happening in 2008, coinciding with
that year's financial crisis.

\subsubsection{Is the variance
constant?}\label{is-the-variance-constant-1}

Let us compute the boxplot for 2-year period in the time series.

\begin{Shaded}
\begin{Highlighting}[]
\CommentTok{\# 1. Get the year for every observation in the TS}
\NormalTok{years }\OtherTok{\textless{}{-}} \FunctionTok{floor}\NormalTok{(}\FunctionTok{time}\NormalTok{(gdpusavar))}

\CommentTok{\# 2. Create the bins (ensure this is the same length as the data)}
\NormalTok{group\_index }\OtherTok{\textless{}{-}} \DecValTok{2} \SpecialCharTok{*}\NormalTok{ (years }\SpecialCharTok{\%/\%} \DecValTok{2}\NormalTok{)}
\NormalTok{group\_labels }\OtherTok{\textless{}{-}} \FunctionTok{paste0}\NormalTok{(group\_index, }\StringTok{"{-}"}\NormalTok{, group\_index }\SpecialCharTok{+} \DecValTok{1}\NormalTok{)}

\CommentTok{\# 3. Use as.numeric() on the TS to prevent the length mismatch error}
\FunctionTok{boxplot}\NormalTok{(}\FunctionTok{as.numeric}\NormalTok{(gdpusavar) }\SpecialCharTok{\textasciitilde{}}\NormalTok{ group\_labels, }
        \AttributeTok{main =} \StringTok{"Quarterly Data by 10{-}Year Intervals"}\NormalTok{,}
        \AttributeTok{xlab =} \StringTok{"Years"}\NormalTok{,}
        \AttributeTok{ylab =} \StringTok{"gdpusavar"}\NormalTok{)}
\end{Highlighting}
\end{Shaded}

\pandocbounded{\includegraphics[keepaspectratio]{CamposCastilloChico_HW1_files/figure-latex/unnamed-chunk-17-1.pdf}}

We notice that the variances differ between periods, even without
considering the 2008 outlier.

The Mean-Variance plot is computed to confirm this.

\begin{Shaded}
\begin{Highlighting}[]
\NormalTok{mean }\OtherTok{=} \FunctionTok{apply}\NormalTok{(}\FunctionTok{matrix}\NormalTok{(gdpusavar, }\AttributeTok{ncol =} \DecValTok{15}\NormalTok{), }\DecValTok{2}\NormalTok{, mean)}
\NormalTok{var }\OtherTok{=} \FunctionTok{apply}\NormalTok{(}\FunctionTok{matrix}\NormalTok{(gdpusavar, }\AttributeTok{ncol =} \DecValTok{15}\NormalTok{), }\DecValTok{2}\NormalTok{, var)}
\FunctionTok{plot}\NormalTok{(var }\SpecialCharTok{\textasciitilde{}}\NormalTok{ mean, }\AttributeTok{main =} \StringTok{"Mean{-}Variance plot"}\NormalTok{)}
\FunctionTok{abline}\NormalTok{(}\FunctionTok{lm}\NormalTok{(var }\SpecialCharTok{\textasciitilde{}}\NormalTok{ mean))}
\end{Highlighting}
\end{Shaded}

\pandocbounded{\includegraphics[keepaspectratio]{CamposCastilloChico_HW1_files/figure-latex/unnamed-chunk-18-1.pdf}}

Indeed, the variance decreases as the mean increases. To fix this, a
logarithmic transformation is applied. Moreover, a constant is added so
as to avoid negative values, which would lead to NaNs.

\begin{Shaded}
\begin{Highlighting}[]
\NormalTok{c }\OtherTok{\textless{}{-}} \FunctionTok{abs}\NormalTok{(}\FunctionTok{min}\NormalTok{(gdpusavar)) }\SpecialCharTok{+} \DecValTok{1}
\NormalTok{lngdpvar }\OtherTok{\textless{}{-}} \FunctionTok{log}\NormalTok{(gdpusavar }\SpecialCharTok{+}\NormalTok{ c)}
\FunctionTok{plot}\NormalTok{(lngdpvar)}
\end{Highlighting}
\end{Shaded}

\pandocbounded{\includegraphics[keepaspectratio]{CamposCastilloChico_HW1_files/figure-latex/unnamed-chunk-19-1.pdf}}

The previous plot shows that the variance has been largely stabilized,
except for the prominent outlier corresponding to the 2008 financial
crisis. We recompute the boxplots and the Mean-Variance plot to ensure
this.

\begin{Shaded}
\begin{Highlighting}[]
\CommentTok{\# 1. Get the year for every observation in the TS}
\NormalTok{years }\OtherTok{\textless{}{-}} \FunctionTok{floor}\NormalTok{(}\FunctionTok{time}\NormalTok{(lngdpvar))}

\CommentTok{\# 2. Create the bins (ensure this is the same length as the data)}
\NormalTok{group\_index }\OtherTok{\textless{}{-}} \DecValTok{2} \SpecialCharTok{*}\NormalTok{ (years }\SpecialCharTok{\%/\%} \DecValTok{2}\NormalTok{)}
\NormalTok{group\_labels }\OtherTok{\textless{}{-}} \FunctionTok{paste0}\NormalTok{(group\_index, }\StringTok{"{-}"}\NormalTok{, group\_index }\SpecialCharTok{+} \DecValTok{1}\NormalTok{)}

\CommentTok{\# 3. Use as.numeric() on the TS to prevent the length mismatch error}
\FunctionTok{boxplot}\NormalTok{(}\FunctionTok{as.numeric}\NormalTok{(lngdpvar) }\SpecialCharTok{\textasciitilde{}}\NormalTok{ group\_labels, }
        \AttributeTok{main =} \StringTok{"Quarterly Data by 2{-}Year Intervals"}\NormalTok{,}
        \AttributeTok{xlab =} \StringTok{"Years"}\NormalTok{,}
        \AttributeTok{ylab =} \StringTok{"ln gdpusavar + c"}\NormalTok{)}
\end{Highlighting}
\end{Shaded}

\pandocbounded{\includegraphics[keepaspectratio]{CamposCastilloChico_HW1_files/figure-latex/unnamed-chunk-20-1.pdf}}

\begin{Shaded}
\begin{Highlighting}[]
\NormalTok{mean }\OtherTok{=} \FunctionTok{apply}\NormalTok{(}\FunctionTok{matrix}\NormalTok{(lngdpvar, }\AttributeTok{ncol =} \DecValTok{15}\NormalTok{), }\DecValTok{2}\NormalTok{, mean)}
\NormalTok{var }\OtherTok{=} \FunctionTok{apply}\NormalTok{(}\FunctionTok{matrix}\NormalTok{(lngdpvar, }\AttributeTok{ncol =} \DecValTok{15}\NormalTok{), }\DecValTok{2}\NormalTok{, var)}
\FunctionTok{plot}\NormalTok{(var }\SpecialCharTok{\textasciitilde{}}\NormalTok{ mean, }\AttributeTok{main =} \StringTok{"Mean{-}Variance plot"}\NormalTok{)}
\FunctionTok{abline}\NormalTok{(}\FunctionTok{lm}\NormalTok{(var }\SpecialCharTok{\textasciitilde{}}\NormalTok{ mean))}
\end{Highlighting}
\end{Shaded}

\pandocbounded{\includegraphics[keepaspectratio]{CamposCastilloChico_HW1_files/figure-latex/unnamed-chunk-21-1.pdf}}

\subsubsection{Is seasonality present?}\label{is-seasonality-present-1}

No, as the time series has already been seasonally adjusted.

\subsubsection{Is the mean constant}\label{is-the-mean-constant-1}

It is not obvious from the boxplot that the mean is constant so we will
look at the ACF and PACF graphs

\begin{Shaded}
\begin{Highlighting}[]
\FunctionTok{par}\NormalTok{(}\AttributeTok{mfrow =} \FunctionTok{c}\NormalTok{(}\DecValTok{1}\NormalTok{, }\DecValTok{2}\NormalTok{))}
\FunctionTok{acf}\NormalTok{(lngdpvar, }\AttributeTok{ylim =} \FunctionTok{c}\NormalTok{(}\SpecialCharTok{{-}}\DecValTok{1}\NormalTok{, }\DecValTok{1}\NormalTok{), }\AttributeTok{lag.max =} \DecValTok{60}\NormalTok{, }\AttributeTok{lwd =} \DecValTok{2}\NormalTok{)}
\FunctionTok{pacf}\NormalTok{(lngdpvar, }\AttributeTok{ylim =} \FunctionTok{c}\NormalTok{(}\SpecialCharTok{{-}}\DecValTok{1}\NormalTok{, }\DecValTok{1}\NormalTok{), }\AttributeTok{lag.max =} \DecValTok{60}\NormalTok{, }\AttributeTok{lwd =} \DecValTok{2}\NormalTok{)}
\end{Highlighting}
\end{Shaded}

\pandocbounded{\includegraphics[keepaspectratio]{CamposCastilloChico_HW1_files/figure-latex/unnamed-chunk-22-1.pdf}}

This already suggests stationarity but let's test if the variance will
change by doing the first difference.

\begin{Shaded}
\begin{Highlighting}[]
\FunctionTok{formatC}\NormalTok{(}\FunctionTok{var}\NormalTok{(lngdpvar), }\AttributeTok{format =} \StringTok{"f"}\NormalTok{, }\AttributeTok{digits =} \DecValTok{7}\NormalTok{)}
\end{Highlighting}
\end{Shaded}

\begin{verbatim}
##    V1         
## V1 "0.0894027"
\end{verbatim}

\begin{Shaded}
\begin{Highlighting}[]
\FunctionTok{formatC}\NormalTok{(}\FunctionTok{var}\NormalTok{(}\FunctionTok{diff}\NormalTok{(lngdpvar)), }\AttributeTok{format =} \StringTok{"f"}\NormalTok{, }\AttributeTok{digits =} \DecValTok{7}\NormalTok{)}
\end{Highlighting}
\end{Shaded}

\begin{verbatim}
##    V1         
## V1 "0.1066934"
\end{verbatim}

The variance increased therefore if we do first difference, it will be
over differentiation.

\textbf{The transformed time series is stationary.}

\begin{Shaded}
\begin{Highlighting}[]
\FunctionTok{par}\NormalTok{(}\AttributeTok{mfrow =} \FunctionTok{c}\NormalTok{(}\DecValTok{1}\NormalTok{, }\DecValTok{2}\NormalTok{))}
\FunctionTok{plot}\NormalTok{(gdpusavar, }\AttributeTok{main =} \StringTok{"X\_t"}\NormalTok{)}
\FunctionTok{plot}\NormalTok{(lngdpvar, }\AttributeTok{main =} \StringTok{"W\_t"}\NormalTok{)}
\FunctionTok{abline}\NormalTok{(}\AttributeTok{h =} \DecValTok{0}\NormalTok{)}
\end{Highlighting}
\end{Shaded}

\pandocbounded{\includegraphics[keepaspectratio]{CamposCastilloChico_HW1_files/figure-latex/unnamed-chunk-24-1.pdf}}

\section{3. Monthly Average Exchange Rate of the
Euro}\label{monthly-average-exchange-rate-of-the-euro}

\subsubsection{Data Loading and
Verification}\label{data-loading-and-verification}

First we load the data by indicating the starting year and setting freq
= 12 to denote monthly data. We view the loaded data in ts format and
verify its attributes.

\begin{Shaded}
\begin{Highlighting}[]
\NormalTok{eurodol }\OtherTok{\textless{}{-}} \FunctionTok{ts}\NormalTok{(}\FunctionTok{read.table}\NormalTok{(}\StringTok{"eurodol.dat"}\NormalTok{, }\AttributeTok{header =}\NormalTok{ F), }\AttributeTok{start =} \DecValTok{1995}\NormalTok{, }\AttributeTok{freq =} \DecValTok{12}\NormalTok{)}
\NormalTok{eurodol}
\end{Highlighting}
\end{Shaded}

\begin{verbatim}
##        Jan   Feb   Mar   Apr   May   Jun   Jul   Aug   Sep   Oct   Nov   Dec
## 1995 1.241 1.259 1.317 1.341 1.322 1.331 1.345 1.304 1.289 1.322 1.324 1.304
## 1996 1.292 1.289 1.281 1.264 1.247 1.253 1.271 1.283 1.269 1.258 1.277 1.250
## 1997 1.216 1.166 1.150 1.145 1.149 1.137 1.105 1.073 1.100 1.120 1.139 1.112
## 1998 1.088 1.088 1.084 1.091 1.109 1.101 1.097 1.102 1.154 1.194 1.164 1.172
## 1999 1.161 1.121 1.088 1.070 1.063 1.038 1.035 1.060 1.050 1.071 1.034 1.011
## 2000 1.014 0.983 0.964 0.947 0.906 0.949 0.940 0.904 0.872 0.855 0.856 0.897
## 2001 0.938 0.922 0.910 0.892 0.874 0.853 0.861 0.900 0.911 0.906 0.888 0.892
## 2002 0.883 0.870 0.876 0.886 0.917 0.955 0.992 0.978 0.981 0.981 1.001 1.018
## 2003 1.062 1.077 1.081 1.085 1.158 1.166 1.137 1.114 1.122 1.169 1.170 1.229
## 2004 1.261 1.265 1.226 1.199 1.201 1.214 1.227 1.218 1.222 1.249 1.299 1.341
## 2005 1.312 1.301 1.320 1.294 1.269 1.216 1.204 1.229 1.226 1.201 1.179 1.186
## 2006 1.210 1.194 1.202 1.227 1.277 1.265 1.268 1.281 1.273 1.261 1.288 1.321
## 2007 1.300 1.307 1.324 1.352 1.351 1.342 1.372 1.362 1.390 1.423 1.468 1.457
## 2008 1.472 1.475 1.553 1.575 1.556 1.555 1.577 1.498 1.437 1.332 1.273 1.345
## 2009 1.324 1.278 1.305 1.319 1.365 1.402 1.409 1.427 1.456 1.482 1.491 1.461
## 2010 1.427 1.369 1.357 1.341 1.257 1.221 1.277 1.289 1.307 1.390 1.366 1.322
## 2011 1.336 1.365 1.400 1.444 1.435 1.439 1.426 1.434 1.377 1.371 1.356 1.318
## 2012 1.290 1.322 1.320 1.316 1.279 1.253 1.229 1.240 1.286 1.297 1.283 1.312
## 2013 1.329 1.336 1.296 1.303 1.298 1.319 1.308 1.331 1.335 1.363 1.349 1.370
## 2014 1.361 1.366 1.382 1.381 1.373 1.359 1.354 1.332 1.290 1.267 1.247 1.233
## 2015 1.162 1.135 1.084 1.078 1.115 1.121 1.100 1.114 1.122 1.124 1.074 1.088
## 2016 1.086 1.109 1.110 1.134 1.131 1.123 1.107 1.121 1.121 1.103 1.080 1.054
## 2017 1.061 1.064 1.068 1.072 1.106 1.123 1.151 1.181 1.191 1.176 1.174 1.184
## 2018 1.220 1.235 1.234 1.228 1.181 1.168 1.169 1.155 1.166 1.148 1.137 1.138
## 2019 1.142 1.135 1.130 1.124 1.118 1.129 1.122 1.113 1.100 1.105 1.105 1.111
\end{verbatim}

\begin{Shaded}
\begin{Highlighting}[]
\FunctionTok{str}\NormalTok{(eurodol)}
\end{Highlighting}
\end{Shaded}

\begin{verbatim}
##  Time-Series [1:300, 1] from 1995 to 2020: 1.24 1.26 1.32 1.34 1.32 ...
##  - attr(*, "dimnames")=List of 2
##   ..$ : NULL
##   ..$ : chr "V1"
\end{verbatim}

\subsubsection{Initial Visualization}\label{initial-visualization}

Now, we plot the data and visualize the time series. I've included lines
to denote the different years and the average of the time series as
well.

\begin{Shaded}
\begin{Highlighting}[]
\FunctionTok{plot}\NormalTok{(eurodol,}
\AttributeTok{main =} \StringTok{"Monthly Average Exchange Rate of the Euro against the Dollar"}\NormalTok{,}
\AttributeTok{ylim =} \FunctionTok{c}\NormalTok{(}\DecValTok{0}\NormalTok{, }\FunctionTok{max}\NormalTok{(eurodol) }\SpecialCharTok{+} \FloatTok{0.1}\NormalTok{),}
\AttributeTok{ylab =} \StringTok{"Average Exchange Rate"}\NormalTok{)}
\FunctionTok{abline}\NormalTok{(}\AttributeTok{v =} \DecValTok{1995}\SpecialCharTok{:}\DecValTok{2020}\NormalTok{, }\AttributeTok{col =} \DecValTok{4}\NormalTok{, }\AttributeTok{lty =} \DecValTok{3}\NormalTok{)}
\FunctionTok{abline}\NormalTok{(}\AttributeTok{h =} \FunctionTok{mean}\NormalTok{(eurodol), }\AttributeTok{col =} \DecValTok{2}\NormalTok{)}
\end{Highlighting}
\end{Shaded}

\pandocbounded{\includegraphics[keepaspectratio]{CamposCastilloChico_HW1_files/figure-latex/plot-euro-1.pdf}}

We see from the graph immediately a non-monotonic trend due to crashes
from the late 90s to early 00s and a spike right before the 2010s.
Non-constant variance is immediately noticeable. There are a few
indicators of yearly seasonality where peaks are generally experienced
in the middle of the year and drop or spike going into the next year.

\subsubsection{Variance Assessment}\label{variance-assessment}

Now we formally assess the variance with the boxplots and mean-variance
plots.

\begin{Shaded}
\begin{Highlighting}[]
\CommentTok{\# Boxplot}
\FunctionTok{boxplot}\NormalTok{(eurodol }\SpecialCharTok{\textasciitilde{}} \FunctionTok{floor}\NormalTok{(}\FunctionTok{time}\NormalTok{(eurodol)), }\AttributeTok{main=}\StringTok{"Boxplot Original"}\NormalTok{)}
\end{Highlighting}
\end{Shaded}

\pandocbounded{\includegraphics[keepaspectratio]{CamposCastilloChico_HW1_files/figure-latex/variance-assessment-euro-1.pdf}}

\begin{Shaded}
\begin{Highlighting}[]
\CommentTok{\# Mean{-}Variance Plot}
\NormalTok{mn }\OtherTok{\textless{}{-}} \FunctionTok{apply}\NormalTok{(}\FunctionTok{matrix}\NormalTok{(eurodol, }\AttributeTok{ncol=}\DecValTok{20}\NormalTok{), }\DecValTok{2}\NormalTok{, mean)}
\NormalTok{var\_val }\OtherTok{\textless{}{-}} \FunctionTok{apply}\NormalTok{(}\FunctionTok{matrix}\NormalTok{(eurodol, }\AttributeTok{ncol=}\DecValTok{20}\NormalTok{), }\DecValTok{2}\NormalTok{, var)}
\FunctionTok{plot}\NormalTok{(var\_val }\SpecialCharTok{\textasciitilde{}}\NormalTok{ mn, }\AttributeTok{main=}\StringTok{"Mean{-}Variance Plot"}\NormalTok{)}
\FunctionTok{abline}\NormalTok{(}\FunctionTok{lm}\NormalTok{(var\_val }\SpecialCharTok{\textasciitilde{}}\NormalTok{ mn))}
\end{Highlighting}
\end{Shaded}

\pandocbounded{\includegraphics[keepaspectratio]{CamposCastilloChico_HW1_files/figure-latex/variance-assessment-euro-2.pdf}}

We can clearly observe non-constant mean. Since we are dealing with
financial data, we transform with the log-transform and compare graphs
again.

\begin{Shaded}
\begin{Highlighting}[]
\NormalTok{lneurodol }\OtherTok{\textless{}{-}} \FunctionTok{log}\NormalTok{(eurodol)}
\FunctionTok{par}\NormalTok{(}\AttributeTok{mfrow =} \FunctionTok{c}\NormalTok{(}\DecValTok{2}\NormalTok{, }\DecValTok{2}\NormalTok{))}
\CommentTok{\# Boxplots}
\FunctionTok{boxplot}\NormalTok{(eurodol }\SpecialCharTok{\textasciitilde{}} \FunctionTok{floor}\NormalTok{(}\FunctionTok{time}\NormalTok{(eurodol)), }\AttributeTok{main =} \StringTok{"Boxplot (Original)"}\NormalTok{)}
\FunctionTok{boxplot}\NormalTok{(lneurodol }\SpecialCharTok{\textasciitilde{}} \FunctionTok{floor}\NormalTok{(}\FunctionTok{time}\NormalTok{(lneurodol)), }\AttributeTok{main =} \StringTok{"Boxplot (Log Transformed)"}\NormalTok{)}
\CommentTok{\# Mean{-}Variance Plots}
\NormalTok{lnmn }\OtherTok{\textless{}{-}} \FunctionTok{apply}\NormalTok{(}\FunctionTok{matrix}\NormalTok{(lneurodol, }\AttributeTok{ncol=}\DecValTok{20}\NormalTok{), }\DecValTok{2}\NormalTok{, mean)}
\NormalTok{lnvar }\OtherTok{\textless{}{-}} \FunctionTok{apply}\NormalTok{(}\FunctionTok{matrix}\NormalTok{(lneurodol, }\AttributeTok{ncol=}\DecValTok{20}\NormalTok{), }\DecValTok{2}\NormalTok{, var)}
\FunctionTok{plot}\NormalTok{(lnvar }\SpecialCharTok{\textasciitilde{}}\NormalTok{ lnmn, }\AttributeTok{main=}\StringTok{"Mean{-}Variance (Log)"}\NormalTok{)}
\FunctionTok{abline}\NormalTok{(}\FunctionTok{lm}\NormalTok{(lnvar }\SpecialCharTok{\textasciitilde{}}\NormalTok{ lnmn))}
\NormalTok{mn }\OtherTok{\textless{}{-}} \FunctionTok{apply}\NormalTok{(}\FunctionTok{matrix}\NormalTok{(eurodol, }\AttributeTok{ncol=}\DecValTok{20}\NormalTok{), }\DecValTok{2}\NormalTok{, mean)}
\NormalTok{var\_val }\OtherTok{\textless{}{-}} \FunctionTok{apply}\NormalTok{(}\FunctionTok{matrix}\NormalTok{(eurodol, }\AttributeTok{ncol=}\DecValTok{20}\NormalTok{), }\DecValTok{2}\NormalTok{, var)}
\FunctionTok{plot}\NormalTok{(var\_val }\SpecialCharTok{\textasciitilde{}}\NormalTok{ mn, }\AttributeTok{main=}\StringTok{"Mean{-}Variance (Original)"}\NormalTok{)}
\FunctionTok{abline}\NormalTok{(}\FunctionTok{lm}\NormalTok{(var\_val }\SpecialCharTok{\textasciitilde{}}\NormalTok{ mn))}
\end{Highlighting}
\end{Shaded}

\pandocbounded{\includegraphics[keepaspectratio]{CamposCastilloChico_HW1_files/figure-latex/log-transform-euro-1.pdf}}

\begin{Shaded}
\begin{Highlighting}[]
\FunctionTok{par}\NormalTok{(}\AttributeTok{mfrow =} \FunctionTok{c}\NormalTok{(}\DecValTok{1}\NormalTok{, }\DecValTok{1}\NormalTok{))}
\end{Highlighting}
\end{Shaded}

\subsubsection{Box-Cox Transformation}\label{box-cox-transformation}

Because the variance is not clearly constant, we double check with a
box-cox transformation.

\begin{Shaded}
\begin{Highlighting}[]
\CommentTok{\# 1. Run the Box{-}Cox function}
\NormalTok{bc }\OtherTok{\textless{}{-}} \FunctionTok{boxcox}\NormalTok{(eurodol }\SpecialCharTok{+} \DecValTok{1} \SpecialCharTok{\textasciitilde{}} \DecValTok{1}\NormalTok{, }\AttributeTok{plotit =} \ConstantTok{TRUE}\NormalTok{)}
\end{Highlighting}
\end{Shaded}

\pandocbounded{\includegraphics[keepaspectratio]{CamposCastilloChico_HW1_files/figure-latex/boxcox-euro-1.pdf}}

\begin{Shaded}
\begin{Highlighting}[]
\CommentTok{\# 2. Extract the exact lambda}
\NormalTok{optimal\_lambda }\OtherTok{\textless{}{-}}\NormalTok{ bc}\SpecialCharTok{$}\NormalTok{x[}\FunctionTok{which.max}\NormalTok{(bc}\SpecialCharTok{$}\NormalTok{y)]}
\NormalTok{optimal\_lambda}
\end{Highlighting}
\end{Shaded}

\begin{verbatim}
## [1] 2
\end{verbatim}

Since the optimal lambda is near 2, we can consider a squared
transformation of our variable and compare it to the log transform.

\begin{Shaded}
\begin{Highlighting}[]
\NormalTok{bceurodol }\OtherTok{\textless{}{-}}\NormalTok{ (eurodol}\SpecialCharTok{\^{}}\NormalTok{optimal\_lambda }\SpecialCharTok{{-}} \DecValTok{1}\NormalTok{) }\SpecialCharTok{/}\NormalTok{ optimal\_lambda}
\FunctionTok{par}\NormalTok{(}\AttributeTok{mfrow =} \FunctionTok{c}\NormalTok{(}\DecValTok{2}\NormalTok{, }\DecValTok{2}\NormalTok{))}
\FunctionTok{boxplot}\NormalTok{(bceurodol }\SpecialCharTok{\textasciitilde{}} \FunctionTok{floor}\NormalTok{(}\FunctionTok{time}\NormalTok{(bceurodol)), }\AttributeTok{main =} \StringTok{"Boxplot (Square)"}\NormalTok{)}
\FunctionTok{boxplot}\NormalTok{(lneurodol }\SpecialCharTok{\textasciitilde{}} \FunctionTok{floor}\NormalTok{(}\FunctionTok{time}\NormalTok{(lneurodol)), }\AttributeTok{main =} \StringTok{"Boxplot (Log)"}\NormalTok{)}
\NormalTok{bcmn }\OtherTok{\textless{}{-}} \FunctionTok{apply}\NormalTok{(}\FunctionTok{matrix}\NormalTok{(bceurodol, }\AttributeTok{ncol=}\DecValTok{20}\NormalTok{), }\DecValTok{2}\NormalTok{, mean)}
\NormalTok{bcvar }\OtherTok{\textless{}{-}} \FunctionTok{apply}\NormalTok{(}\FunctionTok{matrix}\NormalTok{(bceurodol, }\AttributeTok{ncol=}\DecValTok{20}\NormalTok{), }\DecValTok{2}\NormalTok{, var)}
\FunctionTok{plot}\NormalTok{(bcvar }\SpecialCharTok{\textasciitilde{}}\NormalTok{ bcmn, }\AttributeTok{main=}\StringTok{"M{-}V Plot (Square)"}\NormalTok{)}
\FunctionTok{abline}\NormalTok{(}\FunctionTok{lm}\NormalTok{(bcvar }\SpecialCharTok{\textasciitilde{}}\NormalTok{ bcmn))}
\FunctionTok{plot}\NormalTok{(lnvar }\SpecialCharTok{\textasciitilde{}}\NormalTok{ lnmn, }\AttributeTok{main=}\StringTok{"M{-}V Plot (Log)"}\NormalTok{)}
\FunctionTok{abline}\NormalTok{(}\FunctionTok{lm}\NormalTok{(lnvar }\SpecialCharTok{\textasciitilde{}}\NormalTok{ lnmn))}
\end{Highlighting}
\end{Shaded}

\pandocbounded{\includegraphics[keepaspectratio]{CamposCastilloChico_HW1_files/figure-latex/compare-transforms-euro-1.pdf}}

\begin{Shaded}
\begin{Highlighting}[]
\FunctionTok{par}\NormalTok{(}\AttributeTok{mfrow =} \FunctionTok{c}\NormalTok{(}\DecValTok{1}\NormalTok{, }\DecValTok{1}\NormalTok{))}
\end{Highlighting}
\end{Shaded}

\subsubsection{Seasonality Check}\label{seasonality-check}

Now we formally check for seasonality.

\begin{Shaded}
\begin{Highlighting}[]
\FunctionTok{monthplot}\NormalTok{(lneurodol)}
\end{Highlighting}
\end{Shaded}

\pandocbounded{\includegraphics[keepaspectratio]{CamposCastilloChico_HW1_files/figure-latex/seasonality-euro-1.pdf}}

\begin{Shaded}
\begin{Highlighting}[]
\FunctionTok{plot}\NormalTok{(}\FunctionTok{decompose}\NormalTok{(lneurodol))}
\end{Highlighting}
\end{Shaded}

\pandocbounded{\includegraphics[keepaspectratio]{CamposCastilloChico_HW1_files/figure-latex/seasonality-euro-2.pdf}}

\begin{Shaded}
\begin{Highlighting}[]
\FunctionTok{acf}\NormalTok{(lneurodol, }\AttributeTok{ylim =} \FunctionTok{c}\NormalTok{(}\SpecialCharTok{{-}}\DecValTok{1}\NormalTok{,}\DecValTok{1}\NormalTok{), }\AttributeTok{lag.max =} \DecValTok{60}\NormalTok{, }\AttributeTok{col =} \FunctionTok{c}\NormalTok{(}\DecValTok{2}\NormalTok{,}\FunctionTok{rep}\NormalTok{(}\DecValTok{1}\NormalTok{,}\DecValTok{11}\NormalTok{)), }\AttributeTok{lwd=}\DecValTok{2}\NormalTok{, }\AttributeTok{main =} \StringTok{"ACF Plot"}\NormalTok{)}
\end{Highlighting}
\end{Shaded}

\pandocbounded{\includegraphics[keepaspectratio]{CamposCastilloChico_HW1_files/figure-latex/seasonality-euro-3.pdf}}

We see from both the monthplot and the seasonal section of the
decomposition plots that seasonality clearly exists at a yearly level,
hence we adjust for it. The ACF plot reinforces the non-stationarity of
the current series as well.

\subsubsection{Seasonal Differencing}\label{seasonal-differencing}

\begin{Shaded}
\begin{Highlighting}[]
\NormalTok{d12lneurodol }\OtherTok{\textless{}{-}} \FunctionTok{diff}\NormalTok{(lneurodol, }\AttributeTok{lag =} \DecValTok{12}\NormalTok{)}
\FunctionTok{plot}\NormalTok{(d12lneurodol,}
\AttributeTok{main =} \StringTok{"Seasonally Adjusted Log Euro/Dollar"}\NormalTok{,}
\AttributeTok{ylab =} \StringTok{"Log Average Exchange Rate"}\NormalTok{)}
\FunctionTok{abline}\NormalTok{(}\AttributeTok{v =} \DecValTok{1995}\SpecialCharTok{:}\DecValTok{2020}\NormalTok{, }\AttributeTok{col =} \DecValTok{4}\NormalTok{, }\AttributeTok{lty =} \DecValTok{3}\NormalTok{)}
\FunctionTok{abline}\NormalTok{(}\AttributeTok{h =} \FunctionTok{mean}\NormalTok{(d12lneurodol), }\AttributeTok{col =} \DecValTok{2}\NormalTok{)}
\end{Highlighting}
\end{Shaded}

\pandocbounded{\includegraphics[keepaspectratio]{CamposCastilloChico_HW1_files/figure-latex/seasonal-diff-euro-1.pdf}}

\subsubsection{Automatic Differencing
Function}\label{automatic-differencing-function}

We apply regular differencing until variance increases.

\begin{Shaded}
\begin{Highlighting}[]
\NormalTok{auto\_diff }\OtherTok{\textless{}{-}} \ControlFlowTok{function}\NormalTok{(x) \{}
\NormalTok{orig\_name }\OtherTok{\textless{}{-}} \FunctionTok{deparse}\NormalTok{(}\FunctionTok{substitute}\NormalTok{(x))}
\NormalTok{best\_series }\OtherTok{\textless{}{-}}\NormalTok{ x}
\NormalTok{best\_var }\OtherTok{\textless{}{-}} \FunctionTok{var}\NormalTok{(best\_series, }\AttributeTok{na.rm =} \ConstantTok{TRUE}\NormalTok{)}
\NormalTok{diff\_count }\OtherTok{\textless{}{-}} \DecValTok{0}
\FunctionTok{cat}\NormalTok{(}\StringTok{"Initial Variance:"}\NormalTok{, best\_var, }\StringTok{"}\SpecialCharTok{\textbackslash{}n}\StringTok{"}\NormalTok{)}
\ControlFlowTok{repeat}\NormalTok{ \{}
\NormalTok{trial\_series }\OtherTok{\textless{}{-}} \FunctionTok{diff}\NormalTok{(best\_series)}
\NormalTok{trial\_var }\OtherTok{\textless{}{-}} \FunctionTok{var}\NormalTok{(trial\_series, }\AttributeTok{na.rm =} \ConstantTok{TRUE}\NormalTok{)}
\ControlFlowTok{if}\NormalTok{ (trial\_var }\SpecialCharTok{\textgreater{}=}\NormalTok{ best\_var) \{}
\ControlFlowTok{if}\NormalTok{ (diff\_count }\SpecialCharTok{==} \DecValTok{0}\NormalTok{) \{}
\FunctionTok{cat}\NormalTok{(}\StringTok{"Original series is optimal.}\SpecialCharTok{\textbackslash{}n}\StringTok{"}\NormalTok{)}
\NormalTok{\} }\ControlFlowTok{else}\NormalTok{ \{}
\FunctionTok{cat}\NormalTok{(}\StringTok{"Variance increased to"}\NormalTok{, trial\_var, }\StringTok{". Stopping at"}\NormalTok{, diff\_count, }\StringTok{"diff(s).}\SpecialCharTok{\textbackslash{}n}\StringTok{"}\NormalTok{)}
\NormalTok{\}}
\ControlFlowTok{break}
\NormalTok{\}}
\NormalTok{best\_series }\OtherTok{\textless{}{-}}\NormalTok{ trial\_series}
\NormalTok{best\_var }\OtherTok{\textless{}{-}}\NormalTok{ trial\_var}
\NormalTok{diff\_count }\OtherTok{\textless{}{-}}\NormalTok{ diff\_count }\SpecialCharTok{+} \DecValTok{1}
\FunctionTok{cat}\NormalTok{(}\StringTok{"Difference"}\NormalTok{, diff\_count, }\StringTok{"applied. New Variance:"}\NormalTok{, best\_var, }\StringTok{"}\SpecialCharTok{\textbackslash{}n}\StringTok{"}\NormalTok{)}
\ControlFlowTok{if}\NormalTok{ (diff\_count }\SpecialCharTok{\textgreater{}=} \DecValTok{20}\NormalTok{) }\ControlFlowTok{break}
\NormalTok{\}}
\ControlFlowTok{if}\NormalTok{ (diff\_count }\SpecialCharTok{\textgreater{}} \DecValTok{0}\NormalTok{) \{}
\NormalTok{new\_name }\OtherTok{\textless{}{-}} \FunctionTok{paste0}\NormalTok{(}\FunctionTok{paste}\NormalTok{(}\FunctionTok{rep}\NormalTok{(}\StringTok{"d1"}\NormalTok{, diff\_count), }\AttributeTok{collapse =} \StringTok{""}\NormalTok{), orig\_name)}
\FunctionTok{assign}\NormalTok{(new\_name, best\_series, }\AttributeTok{pos =} \DecValTok{1}\NormalTok{)}
\NormalTok{\}}
\FunctionTok{return}\NormalTok{(}\FunctionTok{invisible}\NormalTok{(best\_series))}
\NormalTok{\}}
\CommentTok{\# Run the function}
\FunctionTok{auto\_diff}\NormalTok{(d12lneurodol)}
\end{Highlighting}
\end{Shaded}

\begin{verbatim}
## Initial Variance: 0.009403374 
## Difference 1 applied. New Variance: 0.0010883 
## Variance increased to 0.00144216 . Stopping at 1 diff(s).
\end{verbatim}

\subsubsection{Final Stationary Series}\label{final-stationary-series}

After applying the function, we obtain the series d1d12lneurodol.

\begin{Shaded}
\begin{Highlighting}[]
\FunctionTok{par}\NormalTok{(}\AttributeTok{mfrow =} \FunctionTok{c}\NormalTok{(}\DecValTok{2}\NormalTok{, }\DecValTok{1}\NormalTok{))}
\FunctionTok{plot}\NormalTok{(d1d12lneurodol,}
\AttributeTok{main =} \StringTok{"Detrended \& Seasonally Adjusted Log Euro/Dollar"}\NormalTok{,}
\AttributeTok{ylab =} \StringTok{"Log Average Exchange Rate"}\NormalTok{)}
\FunctionTok{abline}\NormalTok{(}\AttributeTok{v =} \DecValTok{1995}\SpecialCharTok{:}\DecValTok{2020}\NormalTok{, }\AttributeTok{col =} \DecValTok{4}\NormalTok{, }\AttributeTok{lty =} \DecValTok{3}\NormalTok{)}
\FunctionTok{abline}\NormalTok{(}\AttributeTok{h =} \FunctionTok{mean}\NormalTok{(d1d12lneurodol), }\AttributeTok{col =} \DecValTok{2}\NormalTok{)}
\FunctionTok{acf}\NormalTok{(d1d12lneurodol, }\AttributeTok{ylim =} \FunctionTok{c}\NormalTok{(}\SpecialCharTok{{-}}\DecValTok{1}\NormalTok{,}\DecValTok{1}\NormalTok{), }\AttributeTok{lag.max =} \DecValTok{60}\NormalTok{, }\AttributeTok{col =} \FunctionTok{c}\NormalTok{(}\DecValTok{2}\NormalTok{,}\FunctionTok{rep}\NormalTok{(}\DecValTok{1}\NormalTok{,}\DecValTok{11}\NormalTok{)), }\AttributeTok{lwd=}\DecValTok{2}\NormalTok{, }\AttributeTok{main =} \StringTok{"ACF Plot"}\NormalTok{)}
\end{Highlighting}
\end{Shaded}

\pandocbounded{\includegraphics[keepaspectratio]{CamposCastilloChico_HW1_files/figure-latex/final-plot-euro-1.pdf}}

\begin{Shaded}
\begin{Highlighting}[]
\FunctionTok{par}\NormalTok{(}\AttributeTok{mfrow =} \FunctionTok{c}\NormalTok{(}\DecValTok{1}\NormalTok{, }\DecValTok{1}\NormalTok{))}
\end{Highlighting}
\end{Shaded}

\section{4. Number of Individuals Registered as Unemployed
(INEM)}\label{number-of-individuals-registered-as-unemployed-inem}

First we load the data by indicating the starting year and setting freq
= 12 to denote monthly data. We view the loaded data in ts format and
verify its attributes. We also plot the data and visualize the time
series. I've included lines to denote the different years and the
average of the time series as well.

\subsubsection{Data Loading and Initial
Plot}\label{data-loading-and-initial-plot}

\begin{Shaded}
\begin{Highlighting}[]
\NormalTok{atur }\OtherTok{\textless{}{-}} \FunctionTok{ts}\NormalTok{(}\FunctionTok{read.table}\NormalTok{(}\StringTok{"Atur.dat"}\NormalTok{, }\AttributeTok{header =}\NormalTok{ F), }\AttributeTok{start =} \DecValTok{1996}\NormalTok{, }\AttributeTok{freq =} \DecValTok{12}\NormalTok{)}
\FunctionTok{plot}\NormalTok{(atur,}
\AttributeTok{main =} \StringTok{"Monthly Registered Unemployed (INEM)"}\NormalTok{,}
\AttributeTok{ylab =} \StringTok{"Number of Individuals"}\NormalTok{)}
\FunctionTok{abline}\NormalTok{(}\AttributeTok{v =} \DecValTok{1996}\SpecialCharTok{:}\DecValTok{2019}\NormalTok{, }\AttributeTok{col =} \DecValTok{4}\NormalTok{, }\AttributeTok{lty =} \DecValTok{3}\NormalTok{)}
\FunctionTok{abline}\NormalTok{(}\AttributeTok{h =} \FunctionTok{mean}\NormalTok{(atur), }\AttributeTok{col =} \DecValTok{2}\NormalTok{)}
\end{Highlighting}
\end{Shaded}

\pandocbounded{\includegraphics[keepaspectratio]{CamposCastilloChico_HW1_files/figure-latex/load-atur-1.pdf}}

We see from the graph immediately a trend due to abrupt changes from
2007 to 2010. Non-constant variance is not immediately noticeable. There
are a few indicators of yearly seasonality where dips are generally
experienced in the middle of the year and drop or spike going into the
next year. Now we formally assess the variance with the boxplots and
mean-variance plots.

\subsubsection{Variance Stabilization}\label{variance-stabilization}

\begin{Shaded}
\begin{Highlighting}[]
\FunctionTok{boxplot}\NormalTok{(atur }\SpecialCharTok{\textasciitilde{}} \FunctionTok{floor}\NormalTok{(}\FunctionTok{time}\NormalTok{(atur)), }\AttributeTok{main=}\StringTok{"Boxplot Unemployed"}\NormalTok{)}
\end{Highlighting}
\end{Shaded}

\pandocbounded{\includegraphics[keepaspectratio]{CamposCastilloChico_HW1_files/figure-latex/variance-atur-1.pdf}}

\begin{Shaded}
\begin{Highlighting}[]
\NormalTok{mn }\OtherTok{\textless{}{-}} \FunctionTok{apply}\NormalTok{(}\FunctionTok{matrix}\NormalTok{(atur, }\AttributeTok{ncol=}\DecValTok{20}\NormalTok{), }\DecValTok{2}\NormalTok{, mean)}
\NormalTok{var\_val }\OtherTok{\textless{}{-}} \FunctionTok{apply}\NormalTok{(}\FunctionTok{matrix}\NormalTok{(atur, }\AttributeTok{ncol=}\DecValTok{20}\NormalTok{), }\DecValTok{2}\NormalTok{, var)}
\FunctionTok{plot}\NormalTok{(var\_val }\SpecialCharTok{\textasciitilde{}}\NormalTok{ mn, }\AttributeTok{main=}\StringTok{"Mean{-}Variance Plot"}\NormalTok{)}
\FunctionTok{abline}\NormalTok{(}\FunctionTok{lm}\NormalTok{(var\_val }\SpecialCharTok{\textasciitilde{}}\NormalTok{ mn))}
\end{Highlighting}
\end{Shaded}

\pandocbounded{\includegraphics[keepaspectratio]{CamposCastilloChico_HW1_files/figure-latex/variance-atur-2.pdf}}

For the most part, we can see that variance is the similar for most
years except in the period from 2007 to 2010 where variance jumps (as
seen with the long IQRs and outliers in the mean-variance plot). We
apply log transformation to stabilize the variance and lessen the impact
of shocks.

\begin{Shaded}
\begin{Highlighting}[]
\NormalTok{lnatur }\OtherTok{\textless{}{-}} \FunctionTok{log}\NormalTok{(atur)}
\FunctionTok{par}\NormalTok{(}\AttributeTok{mfrow =} \FunctionTok{c}\NormalTok{(}\DecValTok{2}\NormalTok{, }\DecValTok{2}\NormalTok{))}
\FunctionTok{boxplot}\NormalTok{(atur }\SpecialCharTok{\textasciitilde{}} \FunctionTok{floor}\NormalTok{(}\FunctionTok{time}\NormalTok{(atur)), }\AttributeTok{main =} \StringTok{"Original"}\NormalTok{)}
\FunctionTok{boxplot}\NormalTok{(lnatur }\SpecialCharTok{\textasciitilde{}} \FunctionTok{floor}\NormalTok{(}\FunctionTok{time}\NormalTok{(lnatur)), }\AttributeTok{main =} \StringTok{"Log Transformed"}\NormalTok{)}
\NormalTok{lnmn }\OtherTok{\textless{}{-}} \FunctionTok{apply}\NormalTok{(}\FunctionTok{matrix}\NormalTok{(lnatur, }\AttributeTok{ncol=}\DecValTok{20}\NormalTok{), }\DecValTok{2}\NormalTok{, mean)}
\NormalTok{lnvar }\OtherTok{\textless{}{-}} \FunctionTok{apply}\NormalTok{(}\FunctionTok{matrix}\NormalTok{(lnatur, }\AttributeTok{ncol=}\DecValTok{20}\NormalTok{), }\DecValTok{2}\NormalTok{, var)}
\FunctionTok{plot}\NormalTok{(lnvar }\SpecialCharTok{\textasciitilde{}}\NormalTok{ lnmn, }\AttributeTok{main=}\StringTok{"M{-}V Log"}\NormalTok{)}
\FunctionTok{abline}\NormalTok{(}\FunctionTok{lm}\NormalTok{(lnvar }\SpecialCharTok{\textasciitilde{}}\NormalTok{ lnmn))}
\FunctionTok{plot}\NormalTok{(var\_val }\SpecialCharTok{\textasciitilde{}}\NormalTok{ mn, }\AttributeTok{main=}\StringTok{"M{-}V Original"}\NormalTok{)}
\FunctionTok{abline}\NormalTok{(}\FunctionTok{lm}\NormalTok{(var\_val }\SpecialCharTok{\textasciitilde{}}\NormalTok{ mn))}
\end{Highlighting}
\end{Shaded}

\pandocbounded{\includegraphics[keepaspectratio]{CamposCastilloChico_HW1_files/figure-latex/log-atur-1.pdf}}

\begin{Shaded}
\begin{Highlighting}[]
\FunctionTok{par}\NormalTok{(}\AttributeTok{mfrow =} \FunctionTok{c}\NormalTok{(}\DecValTok{1}\NormalTok{, }\DecValTok{1}\NormalTok{))}
\end{Highlighting}
\end{Shaded}

We now observe, especially in the mean-variance plot, that variance is
now more constant even in the presence of outliers. We now observe for
seasonality and differencing.

\subsubsection{Seasonality and
Differencing}\label{seasonality-and-differencing}

\begin{Shaded}
\begin{Highlighting}[]
\FunctionTok{monthplot}\NormalTok{(lnatur)}
\end{Highlighting}
\end{Shaded}

\pandocbounded{\includegraphics[keepaspectratio]{CamposCastilloChico_HW1_files/figure-latex/seasonality-atur-1.pdf}}

\begin{Shaded}
\begin{Highlighting}[]
\FunctionTok{plot}\NormalTok{(}\FunctionTok{decompose}\NormalTok{(lnatur))}
\end{Highlighting}
\end{Shaded}

\pandocbounded{\includegraphics[keepaspectratio]{CamposCastilloChico_HW1_files/figure-latex/seasonality-atur-2.pdf}}

\begin{Shaded}
\begin{Highlighting}[]
\FunctionTok{acf}\NormalTok{(lnatur, }\AttributeTok{ylim =} \FunctionTok{c}\NormalTok{(}\SpecialCharTok{{-}}\DecValTok{1}\NormalTok{,}\DecValTok{1}\NormalTok{), }\AttributeTok{lag.max =} \DecValTok{60}\NormalTok{, }\AttributeTok{col =} \FunctionTok{c}\NormalTok{(}\DecValTok{2}\NormalTok{,}\FunctionTok{rep}\NormalTok{(}\DecValTok{1}\NormalTok{,}\DecValTok{11}\NormalTok{)), }\AttributeTok{lwd=}\DecValTok{2}\NormalTok{)}
\end{Highlighting}
\end{Shaded}

\pandocbounded{\includegraphics[keepaspectratio]{CamposCastilloChico_HW1_files/figure-latex/seasonality-atur-3.pdf}}

\begin{Shaded}
\begin{Highlighting}[]
\NormalTok{d12lnatur }\OtherTok{\textless{}{-}} \FunctionTok{diff}\NormalTok{(lnatur, }\AttributeTok{lag =} \DecValTok{12}\NormalTok{)}
\FunctionTok{auto\_diff}\NormalTok{(d12lnatur)}
\end{Highlighting}
\end{Shaded}

\begin{verbatim}
## Initial Variance: 0.01306892 
## Difference 1 applied. New Variance: 0.0001692097 
## Variance increased to 0.0002012453 . Stopping at 1 diff(s).
\end{verbatim}

We see from both the monthplot and the seasonal section of the
decomposition plots that seasonality clearly exists at a yearly level,
hence we adjust for it. The ACF plot reinforces the non-stationarity of
the current series as well. We remove seasonality.

From the resulting plot too, a trend is clearly seen as the mean is not
constant especially at the shock. We apply regular differencing until
variance increases. For this, we use a function that applies
differencing until the variance of the series increases. It returns the
previously differenced series when true.

To access the new series, we open out environment and look for a series
containing d1. After applying the function, we obtain the series
d1d12lnatur, meaning we applied one difference.

\subsubsection{Final Stationary Series
(Unemployment)}\label{final-stationary-series-unemployment}

We now see the detrended plot of the time series. We also check the ACF
plot for non-stationarity.

\begin{Shaded}
\begin{Highlighting}[]
\FunctionTok{par}\NormalTok{(}\AttributeTok{mfrow =} \FunctionTok{c}\NormalTok{(}\DecValTok{2}\NormalTok{, }\DecValTok{1}\NormalTok{))}
\FunctionTok{plot}\NormalTok{(d1d12lnatur,}
\AttributeTok{main =} \StringTok{"Final Stationary Unemployment Series"}\NormalTok{,}
\AttributeTok{ylab =} \StringTok{"Log Number of Individuals"}\NormalTok{)}
\FunctionTok{abline}\NormalTok{(}\AttributeTok{v =} \DecValTok{1996}\SpecialCharTok{:}\DecValTok{2019}\NormalTok{, }\AttributeTok{col =} \DecValTok{4}\NormalTok{, }\AttributeTok{lty =} \DecValTok{3}\NormalTok{)}
\FunctionTok{abline}\NormalTok{(}\AttributeTok{h =} \FunctionTok{mean}\NormalTok{(d1d12lnatur), }\AttributeTok{col =} \DecValTok{2}\NormalTok{)}
\FunctionTok{acf}\NormalTok{(d1d12lnatur, }\AttributeTok{ylim =} \FunctionTok{c}\NormalTok{(}\SpecialCharTok{{-}}\DecValTok{1}\NormalTok{,}\DecValTok{1}\NormalTok{), }\AttributeTok{lag.max =} \DecValTok{60}\NormalTok{, }\AttributeTok{col =} \FunctionTok{c}\NormalTok{(}\DecValTok{2}\NormalTok{,}\FunctionTok{rep}\NormalTok{(}\DecValTok{1}\NormalTok{,}\DecValTok{11}\NormalTok{)), }\AttributeTok{lwd=}\DecValTok{2}\NormalTok{)}
\end{Highlighting}
\end{Shaded}

\pandocbounded{\includegraphics[keepaspectratio]{CamposCastilloChico_HW1_files/figure-latex/final-plot-atur-1.pdf}}

\begin{Shaded}
\begin{Highlighting}[]
\FunctionTok{par}\NormalTok{(}\AttributeTok{mfrow =} \FunctionTok{c}\NormalTok{(}\DecValTok{1}\NormalTok{, }\DecValTok{1}\NormalTok{))}
\end{Highlighting}
\end{Shaded}

Now we stop as have obtained a stationary time series as seen from both
the time series plot and ACF plot.

\section{5. Number of Monthly International Flight Passengers at El Prat
(BCN)}\label{number-of-monthly-international-flight-passengers-at-el-prat-bcn}

\subsection{Configuration and Data
Loading}\label{configuration-and-data-loading}

\begin{Shaded}
\begin{Highlighting}[]
\CommentTok{\# Set working directory and parameters}
\NormalTok{data\_file }\OtherTok{\textless{}{-}} \StringTok{"AirBCN.dat"}
\NormalTok{series\_name }\OtherTok{\textless{}{-}} \StringTok{"AIRBCN"}

\NormalTok{start\_year }\OtherTok{\textless{}{-}} \DecValTok{1990}
\NormalTok{start\_period }\OtherTok{\textless{}{-}} \DecValTok{1}
\NormalTok{freq }\OtherTok{\textless{}{-}} \DecValTok{12}  
\NormalTok{period\_length }\OtherTok{\textless{}{-}} \DecValTok{12} 
\NormalTok{xlabel }\OtherTok{\textless{}{-}} \StringTok{"month"}
\NormalTok{ylabel }\OtherTok{\textless{}{-}} \StringTok{"passengers"}
\end{Highlighting}
\end{Shaded}

\subsection{1. Load Data as Time Series
Object}\label{load-data-as-time-series-object}

\begin{Shaded}
\begin{Highlighting}[]
\CommentTok{\# Read data}
\NormalTok{data\_raw }\OtherTok{\textless{}{-}} \FunctionTok{read.table}\NormalTok{(data\_file)}
\NormalTok{series }\OtherTok{\textless{}{-}} \FunctionTok{ts}\NormalTok{(data\_raw[,}\DecValTok{1}\NormalTok{], }\AttributeTok{start=}\FunctionTok{c}\NormalTok{(start\_year, start\_period), }\AttributeTok{frequency=}\NormalTok{freq)}

\CommentTok{\# Display properties}
\FunctionTok{cat}\NormalTok{(}\StringTok{"Series Name:"}\NormalTok{, series\_name, }\StringTok{"}\SpecialCharTok{\textbackslash{}n}\StringTok{"}\NormalTok{)}
\end{Highlighting}
\end{Shaded}

\begin{verbatim}
## Series Name: AIRBCN
\end{verbatim}

\begin{Shaded}
\begin{Highlighting}[]
\FunctionTok{cat}\NormalTok{(}\StringTok{"Class:"}\NormalTok{, }\FunctionTok{class}\NormalTok{(series), }\StringTok{"}\SpecialCharTok{\textbackslash{}n}\StringTok{"}\NormalTok{)}
\end{Highlighting}
\end{Shaded}

\begin{verbatim}
## Class: ts
\end{verbatim}

\begin{Shaded}
\begin{Highlighting}[]
\FunctionTok{cat}\NormalTok{(}\StringTok{"Length:"}\NormalTok{, }\FunctionTok{length}\NormalTok{(series), }\StringTok{"observations}\SpecialCharTok{\textbackslash{}n}\StringTok{"}\NormalTok{)}
\end{Highlighting}
\end{Shaded}

\begin{verbatim}
## Length: 360 observations
\end{verbatim}

\begin{Shaded}
\begin{Highlighting}[]
\FunctionTok{cat}\NormalTok{(}\StringTok{"Start:"}\NormalTok{, }\FunctionTok{start}\NormalTok{(series)[}\DecValTok{1}\NormalTok{], }\StringTok{"Period"}\NormalTok{, }\FunctionTok{start}\NormalTok{(series)[}\DecValTok{2}\NormalTok{], }\StringTok{"}\SpecialCharTok{\textbackslash{}n}\StringTok{"}\NormalTok{)}
\end{Highlighting}
\end{Shaded}

\begin{verbatim}
## Start: 1990 Period 1
\end{verbatim}

\begin{Shaded}
\begin{Highlighting}[]
\FunctionTok{cat}\NormalTok{(}\StringTok{"End:"}\NormalTok{, }\FunctionTok{end}\NormalTok{(series)[}\DecValTok{1}\NormalTok{], }\StringTok{"Period"}\NormalTok{, }\FunctionTok{end}\NormalTok{(series)[}\DecValTok{2}\NormalTok{], }\StringTok{"}\SpecialCharTok{\textbackslash{}n}\StringTok{"}\NormalTok{)}
\end{Highlighting}
\end{Shaded}

\begin{verbatim}
## End: 2019 Period 12
\end{verbatim}

\begin{Shaded}
\begin{Highlighting}[]
\FunctionTok{cat}\NormalTok{(}\StringTok{"Frequency:"}\NormalTok{, }\FunctionTok{frequency}\NormalTok{(series), }\StringTok{"}\SpecialCharTok{\textbackslash{}n\textbackslash{}n}\StringTok{"}\NormalTok{)}
\end{Highlighting}
\end{Shaded}

\begin{verbatim}
## Frequency: 12
\end{verbatim}

\begin{Shaded}
\begin{Highlighting}[]
\CommentTok{\# Show first and last values}
\FunctionTok{head}\NormalTok{(series)}
\end{Highlighting}
\end{Shaded}

\begin{verbatim}
## [1] 215140 206775 265051 337802 305500 317929
\end{verbatim}

\begin{Shaded}
\begin{Highlighting}[]
\FunctionTok{tail}\NormalTok{(series)}
\end{Highlighting}
\end{Shaded}

\begin{verbatim}
## [1] 4001787 4132202 3815026 3427111 2624053 2629175
\end{verbatim}

\subsection{2. Generate Graph and Describe Visual
Characteristics}\label{generate-graph-and-describe-visual-characteristics}

The data seems to have an exponential trend, variance increasing with
the level of the series, and some seasonal pattern.

\begin{Shaded}
\begin{Highlighting}[]
\FunctionTok{plot}\NormalTok{(series, }\AttributeTok{type=}\StringTok{"o"}\NormalTok{, }\AttributeTok{main=}\NormalTok{series\_name, }\AttributeTok{ylab=}\NormalTok{ylabel, }\AttributeTok{xlab=}\NormalTok{xlabel)}
\FunctionTok{abline}\NormalTok{(}\AttributeTok{h=}\FunctionTok{mean}\NormalTok{(series), }\AttributeTok{col=}\DecValTok{2}\NormalTok{, }\AttributeTok{lwd=}\DecValTok{2}\NormalTok{)}
\FunctionTok{grid}\NormalTok{()}
\end{Highlighting}
\end{Shaded}

\pandocbounded{\includegraphics[keepaspectratio]{CamposCastilloChico_HW1_files/figure-latex/plot_original-1.pdf}}

\subsection{3. Analyze Variance
Constancy}\label{analyze-variance-constancy}

\subsubsection{3.1 Mean-Variance
Relationship}\label{mean-variance-relationship}

Variance is highly correlated with the mean (0.964).

\begin{Shaded}
\begin{Highlighting}[]
\CommentTok{\# Create matrix for analysis (each column = one period)}
\NormalTok{n\_obs }\OtherTok{\textless{}{-}} \FunctionTok{length}\NormalTok{(series)}
\NormalTok{n\_groups }\OtherTok{\textless{}{-}} \FunctionTok{floor}\NormalTok{(n\_obs }\SpecialCharTok{/}\NormalTok{ period\_length)}
\NormalTok{series\_matrix }\OtherTok{\textless{}{-}} \FunctionTok{matrix}\NormalTok{(series, }\AttributeTok{ncol=}\NormalTok{n\_groups)}

\CommentTok{\# Calculate mean and variance for each period}
\NormalTok{m }\OtherTok{\textless{}{-}} \FunctionTok{apply}\NormalTok{(series\_matrix, }\DecValTok{2}\NormalTok{, mean)}
\NormalTok{v }\OtherTok{\textless{}{-}} \FunctionTok{apply}\NormalTok{(series\_matrix, }\DecValTok{2}\NormalTok{, var)}
\end{Highlighting}
\end{Shaded}

\begin{Shaded}
\begin{Highlighting}[]
\CommentTok{\# Mean{-}variance plot}
\FunctionTok{plot}\NormalTok{(v }\SpecialCharTok{\textasciitilde{}}\NormalTok{ m,}
  \AttributeTok{main=}\FunctionTok{paste}\NormalTok{(}\StringTok{"Mean vs Variance:"}\NormalTok{, series\_name),}
  \AttributeTok{xlab=}\FunctionTok{paste}\NormalTok{(period\_length,}\StringTok{"Period Mean"}\NormalTok{, }\AttributeTok{sep=}\StringTok{"{-}"}\NormalTok{), }
  \AttributeTok{ylab=}\FunctionTok{paste}\NormalTok{(period\_length,}\StringTok{"Period Variance"}\NormalTok{, }\AttributeTok{sep=}\StringTok{"{-}"}\NormalTok{), }
  \AttributeTok{pch=}\DecValTok{19}\NormalTok{, }\AttributeTok{col=}\StringTok{"blue"}\NormalTok{,}\AttributeTok{xaxt=}\StringTok{"n"}\NormalTok{, }\AttributeTok{yaxt=}\StringTok{"n"}\NormalTok{)}
\FunctionTok{axis}\NormalTok{(}\DecValTok{1}\NormalTok{, }\AttributeTok{at=}\FunctionTok{axTicks}\NormalTok{(}\DecValTok{1}\NormalTok{), }\AttributeTok{labels=}\FunctionTok{paste0}\NormalTok{(}\FunctionTok{round}\NormalTok{(}\FunctionTok{axTicks}\NormalTok{(}\DecValTok{1}\NormalTok{)}\SpecialCharTok{/}\FloatTok{1e6}\NormalTok{, }\DecValTok{2}\NormalTok{), }\StringTok{"Mn"}\NormalTok{))}
\FunctionTok{axis}\NormalTok{(}\DecValTok{2}\NormalTok{, }\AttributeTok{at=}\FunctionTok{axTicks}\NormalTok{(}\DecValTok{2}\NormalTok{), }\AttributeTok{labels=}\FunctionTok{paste0}\NormalTok{(}\FunctionTok{round}\NormalTok{(}\FunctionTok{axTicks}\NormalTok{(}\DecValTok{2}\NormalTok{)}\SpecialCharTok{/}\FloatTok{1e9}\NormalTok{, }\DecValTok{2}\NormalTok{), }\StringTok{"Bn"}\NormalTok{))}
\FunctionTok{abline}\NormalTok{(}\FunctionTok{lm}\NormalTok{(v }\SpecialCharTok{\textasciitilde{}}\NormalTok{ m))}
\FunctionTok{box}\NormalTok{()}
\end{Highlighting}
\end{Shaded}

\pandocbounded{\includegraphics[keepaspectratio]{CamposCastilloChico_HW1_files/figure-latex/mean_variance_plot_original-1.pdf}}

\begin{Shaded}
\begin{Highlighting}[]
\CommentTok{\# Correlation}
\FunctionTok{cat}\NormalTok{(}\StringTok{"Correlation between mean and variance:"}\NormalTok{, }\FunctionTok{round}\NormalTok{(}\FunctionTok{cor}\NormalTok{(m, v), }\DecValTok{3}\NormalTok{), }\StringTok{"}\SpecialCharTok{\textbackslash{}n}\StringTok{"}\NormalTok{)}
\end{Highlighting}
\end{Shaded}

\begin{verbatim}
## Correlation between mean and variance: 0.964
\end{verbatim}

\subsubsection{3.2 Boxplots for Different
Periods}\label{boxplots-for-different-periods}

The observations from the mean-variance plot are confirmed by the
boxplots, which show increasing mean and higher box heights over time.

\begin{Shaded}
\begin{Highlighting}[]
\FunctionTok{boxplot}\NormalTok{(series }\SpecialCharTok{\textasciitilde{}} \FunctionTok{floor}\NormalTok{(}\FunctionTok{time}\NormalTok{(series)), }
        \AttributeTok{main=}\FunctionTok{paste}\NormalTok{(}\StringTok{"Yearly Distributions:"}\NormalTok{, series\_name),}
        \AttributeTok{xlab=}\StringTok{"Year"}\NormalTok{, }\AttributeTok{ylab=}\NormalTok{ylabel, }\AttributeTok{yaxt=}\StringTok{"n"}\NormalTok{)}
\FunctionTok{axis}\NormalTok{(}\DecValTok{2}\NormalTok{, }\AttributeTok{at=}\FunctionTok{axTicks}\NormalTok{(}\DecValTok{2}\NormalTok{), }\AttributeTok{labels=}\FunctionTok{paste0}\NormalTok{(}\FunctionTok{round}\NormalTok{(}\FunctionTok{axTicks}\NormalTok{(}\DecValTok{2}\NormalTok{)}\SpecialCharTok{/}\FloatTok{1e6}\NormalTok{, }\DecValTok{2}\NormalTok{), }\StringTok{"Mn"}\NormalTok{))}
\end{Highlighting}
\end{Shaded}

\pandocbounded{\includegraphics[keepaspectratio]{CamposCastilloChico_HW1_files/figure-latex/boxplots_original-1.pdf}}

\subsubsection{3.3 Apply Logarithmic
Transformation}\label{apply-logarithmic-transformation}

The logarithmic transformation seem to stabilize the variance because
the variation across time is more uniform.

\begin{Shaded}
\begin{Highlighting}[]
\NormalTok{lnseries }\OtherTok{\textless{}{-}} \FunctionTok{log}\NormalTok{(series)}
\FunctionTok{plot}\NormalTok{(lnseries, }\AttributeTok{type=}\StringTok{"o"}\NormalTok{, }
  \AttributeTok{main=}\FunctionTok{paste}\NormalTok{(series\_name, }\StringTok{"(ln transformed)"}\NormalTok{),}
  \AttributeTok{ylab=}\FunctionTok{paste}\NormalTok{(}\StringTok{"ln("}\NormalTok{, ylabel, }\StringTok{")"}\NormalTok{, }\AttributeTok{sep=}\StringTok{""}\NormalTok{),}
  \AttributeTok{xlab=}\StringTok{""}\NormalTok{)}
\FunctionTok{abline}\NormalTok{(}\AttributeTok{h=}\FunctionTok{mean}\NormalTok{(lnseries), }\AttributeTok{col=}\DecValTok{2}\NormalTok{, }\AttributeTok{lwd=}\DecValTok{2}\NormalTok{)}
\end{Highlighting}
\end{Shaded}

\pandocbounded{\includegraphics[keepaspectratio]{CamposCastilloChico_HW1_files/figure-latex/log_transform-1.pdf}}

\subsubsection{3.4 Re-assess Variance
Constancy}\label{re-assess-variance-constancy}

Correlation between mean and variance is significantly reduced to 0.307
but still need to check if variance is constant over time.

\begin{Shaded}
\begin{Highlighting}[]
\NormalTok{series\_matrix }\OtherTok{\textless{}{-}} \FunctionTok{matrix}\NormalTok{(lnseries, }\AttributeTok{ncol=}\NormalTok{n\_groups)}
\NormalTok{m }\OtherTok{\textless{}{-}} \FunctionTok{apply}\NormalTok{(series\_matrix, }\DecValTok{2}\NormalTok{, mean)}
\NormalTok{v }\OtherTok{\textless{}{-}} \FunctionTok{apply}\NormalTok{(series\_matrix, }\DecValTok{2}\NormalTok{, var)}

\FunctionTok{plot}\NormalTok{(v }\SpecialCharTok{\textasciitilde{}}\NormalTok{ m, }
  \AttributeTok{main=}\FunctionTok{paste}\NormalTok{(}\StringTok{"Mean vs Variance:"}\NormalTok{, series\_name, }\StringTok{"(ln transformed)"}\NormalTok{),}
  \AttributeTok{xlab=}\FunctionTok{paste}\NormalTok{(period\_length,}\StringTok{"Period Mean"}\NormalTok{, }\AttributeTok{sep=}\StringTok{"{-}"}\NormalTok{), }
  \AttributeTok{ylab=}\FunctionTok{paste}\NormalTok{(period\_length,}\StringTok{"Period Variance"}\NormalTok{, }\AttributeTok{sep=}\StringTok{"{-}"}\NormalTok{), }\AttributeTok{pch=}\DecValTok{19}\NormalTok{, }\AttributeTok{col=}\StringTok{"blue"}\NormalTok{)}
\FunctionTok{abline}\NormalTok{(}\FunctionTok{lm}\NormalTok{(v }\SpecialCharTok{\textasciitilde{}}\NormalTok{ m), }\AttributeTok{col=}\DecValTok{2}\NormalTok{, }\AttributeTok{lwd=}\DecValTok{2}\NormalTok{)}
\end{Highlighting}
\end{Shaded}

\pandocbounded{\includegraphics[keepaspectratio]{CamposCastilloChico_HW1_files/figure-latex/mean_variance_plot_log-1.pdf}}

\begin{Shaded}
\begin{Highlighting}[]
\FunctionTok{cat}\NormalTok{(}\StringTok{"Correlation (ln transformed):"}\NormalTok{, }\FunctionTok{round}\NormalTok{(}\FunctionTok{cor}\NormalTok{(m, v), }\DecValTok{3}\NormalTok{), }\StringTok{"}\SpecialCharTok{\textbackslash{}n}\StringTok{"}\NormalTok{)}
\end{Highlighting}
\end{Shaded}

\begin{verbatim}
## Correlation (ln transformed): 0.307
\end{verbatim}

The boxplots show more consistent variance across years after log
transformation.

\begin{Shaded}
\begin{Highlighting}[]
\FunctionTok{boxplot}\NormalTok{(lnseries }\SpecialCharTok{\textasciitilde{}} \FunctionTok{floor}\NormalTok{(}\FunctionTok{time}\NormalTok{(lnseries)), }
        \AttributeTok{main=}\FunctionTok{paste}\NormalTok{(}\StringTok{"Boxplots by Year:"}\NormalTok{, series\_name, }\StringTok{"(ln transformed)"}\NormalTok{),}
        \AttributeTok{xlab=}\StringTok{"Year"}\NormalTok{, }\AttributeTok{ylab=}\FunctionTok{paste}\NormalTok{(}\StringTok{"ln("}\NormalTok{, ylabel, }\StringTok{")"}\NormalTok{, }\AttributeTok{sep=}\StringTok{""}\NormalTok{))}
\end{Highlighting}
\end{Shaded}

\pandocbounded{\includegraphics[keepaspectratio]{CamposCastilloChico_HW1_files/figure-latex/boxplots_log-1.pdf}}

The ln transformation appears to stabilize the variance. To achieve weak
stationarity, the remaining task is to stabilize the mean and ensure the
autocovariance depends only on lag. We do this by examining seasonal
patterns and applying differencing as needed.

\begin{Shaded}
\begin{Highlighting}[]
\CommentTok{\# Update working series}
\NormalTok{working\_series\_name }\OtherTok{\textless{}{-}} \FunctionTok{paste}\NormalTok{(}\StringTok{"ln("}\NormalTok{, series\_name, }\StringTok{")"}\NormalTok{, }\AttributeTok{sep=}\StringTok{""}\NormalTok{)}
\NormalTok{working\_ylabel }\OtherTok{\textless{}{-}} \FunctionTok{paste}\NormalTok{(}\StringTok{"ln("}\NormalTok{, ylabel, }\StringTok{")"}\NormalTok{, }\AttributeTok{sep=}\StringTok{""}\NormalTok{)}
\NormalTok{working\_series }\OtherTok{\textless{}{-}}\NormalTok{ lnseries}
\end{Highlighting}
\end{Shaded}

\subsection{4. Examine Seasonal
Patterns}\label{examine-seasonal-patterns}

\subsubsection{4.1 Time Series Plot (Zoom)}\label{time-series-plot-zoom}

The series shows a seasonal pattern with peak at the middle of each year
and troughs at the start and end of each year.

\begin{Shaded}
\begin{Highlighting}[]
\CommentTok{\# Plot a few years to visualize seasonality}
\NormalTok{start\_zoom }\OtherTok{\textless{}{-}} \FunctionTok{start}\NormalTok{(working\_series)[}\DecValTok{1}\NormalTok{]}
\NormalTok{end\_zoom }\OtherTok{\textless{}{-}} \FunctionTok{min}\NormalTok{(start\_zoom }\SpecialCharTok{+} \DecValTok{3}\NormalTok{, }\FunctionTok{end}\NormalTok{(working\_series)[}\DecValTok{1}\NormalTok{])}

\FunctionTok{plot}\NormalTok{(}\FunctionTok{window}\NormalTok{(working\_series, }\AttributeTok{start=}\NormalTok{start\_zoom, }\AttributeTok{end=}\NormalTok{end\_zoom), }
     \AttributeTok{type=}\StringTok{"o"}\NormalTok{, }\AttributeTok{main=}\FunctionTok{paste}\NormalTok{(working\_series\_name),}
     \AttributeTok{ylab=}\NormalTok{working\_ylabel)}
\FunctionTok{abline}\NormalTok{(}\AttributeTok{v=}\NormalTok{start\_zoom}\SpecialCharTok{:}\NormalTok{end\_zoom, }\AttributeTok{col=}\DecValTok{4}\NormalTok{, }\AttributeTok{lty=}\DecValTok{3}\NormalTok{)}
\end{Highlighting}
\end{Shaded}

\pandocbounded{\includegraphics[keepaspectratio]{CamposCastilloChico_HW1_files/figure-latex/seasonal_plot-1.pdf}}

The monthly means show the seasonal patter, with the highest in August
and the lowest in January, February, November, and December. Note that
there is still an upward for all months over the years.

\begin{Shaded}
\begin{Highlighting}[]
\CommentTok{\# Monthplot }
\FunctionTok{monthplot}\NormalTok{(working\_series, }\AttributeTok{main=}\FunctionTok{paste}\NormalTok{(}\StringTok{"Month Plot:"}\NormalTok{, series\_name),}
            \AttributeTok{ylab=}\NormalTok{ylabel, }\AttributeTok{xlab=}\StringTok{"Month"}\NormalTok{)}
\end{Highlighting}
\end{Shaded}

\pandocbounded{\includegraphics[keepaspectratio]{CamposCastilloChico_HW1_files/figure-latex/monthly_means-1.pdf}}

\subsubsection{4.3 Autocorrelation Function
(ACF)}\label{autocorrelation-function-acf}

ACF has slow decay and the 12-month seasonal cycle is more obvious. PACF
confirms the 12-lag seasonality.

\begin{Shaded}
\begin{Highlighting}[]
\FunctionTok{par}\NormalTok{(}\AttributeTok{mfrow=}\FunctionTok{c}\NormalTok{(}\DecValTok{1}\NormalTok{,}\DecValTok{2}\NormalTok{))}
\NormalTok{lag\_colors }\OtherTok{\textless{}{-}} \FunctionTok{c}\NormalTok{(}\DecValTok{2}\NormalTok{, }\FunctionTok{rep}\NormalTok{(}\DecValTok{1}\NormalTok{, }\DecValTok{11}\NormalTok{))  }\CommentTok{\# Highlight lag 12}
\FunctionTok{acf}\NormalTok{(working\_series, }\AttributeTok{ylim=}\FunctionTok{c}\NormalTok{(}\SpecialCharTok{{-}}\DecValTok{1}\NormalTok{,}\DecValTok{1}\NormalTok{), }\AttributeTok{lag.max=}\FunctionTok{max}\NormalTok{(}\DecValTok{60}\NormalTok{, }\DecValTok{3}\SpecialCharTok{*}\NormalTok{freq),}
    \AttributeTok{col=}\NormalTok{lag\_colors, }\AttributeTok{lwd=}\DecValTok{2}\NormalTok{, }\AttributeTok{main=}\FunctionTok{paste}\NormalTok{(}\StringTok{"ACF:"}\NormalTok{, working\_series\_name))}

\FunctionTok{pacf}\NormalTok{(working\_series, }\AttributeTok{ylim=}\FunctionTok{c}\NormalTok{(}\SpecialCharTok{{-}}\DecValTok{1}\NormalTok{,}\DecValTok{1}\NormalTok{), }\AttributeTok{lag.max=}\FunctionTok{max}\NormalTok{(}\DecValTok{60}\NormalTok{, }\DecValTok{3}\SpecialCharTok{*}\NormalTok{freq), }
    \AttributeTok{col=}\NormalTok{lag\_colors, }\AttributeTok{lwd=}\DecValTok{2}\NormalTok{, }\AttributeTok{main=}\FunctionTok{paste}\NormalTok{(}\StringTok{"PACF:"}\NormalTok{, working\_series\_name)}
\NormalTok{)}
\end{Highlighting}
\end{Shaded}

\pandocbounded{\includegraphics[keepaspectratio]{CamposCastilloChico_HW1_files/figure-latex/acf-1.pdf}}

\subsubsection{4.4 Seasonal Differencing}\label{seasonal-differencing-1}

We apply seasonal differencing with lag 12 to remove the seasonal
component.

\begin{Shaded}
\begin{Highlighting}[]
\NormalTok{lag }\OtherTok{=} \DecValTok{12}
\NormalTok{acf\_values }\OtherTok{\textless{}{-}} \FunctionTok{acf}\NormalTok{(working\_series, }\AttributeTok{plot=}\ConstantTok{FALSE}\NormalTok{, }\AttributeTok{lag.max=}\NormalTok{freq)}\SpecialCharTok{$}\NormalTok{acf}
\NormalTok{d12series }\OtherTok{\textless{}{-}} \FunctionTok{diff}\NormalTok{(working\_series, }\AttributeTok{lag=}\NormalTok{lag)}
\end{Highlighting}
\end{Shaded}

\begin{Shaded}
\begin{Highlighting}[]
\NormalTok{working\_series\_name }\OtherTok{\textless{}{-}} \FunctionTok{paste0}\NormalTok{(}\StringTok{"Δ"}\NormalTok{, lag, }\StringTok{" ln("}\NormalTok{, series\_name, }\StringTok{")"}\NormalTok{, }\AttributeTok{sep=}\StringTok{""}\NormalTok{)}
\NormalTok{working\_ylabel }\OtherTok{\textless{}{-}} \FunctionTok{paste0}\NormalTok{(}\StringTok{"Δ"}\NormalTok{, lag, }\StringTok{" ln("}\NormalTok{, ylabel, }\StringTok{")"}\NormalTok{, }\AttributeTok{sep=}\StringTok{""}\NormalTok{)}
\NormalTok{working\_series }\OtherTok{\textless{}{-}}\NormalTok{ d12series}
\end{Highlighting}
\end{Shaded}

\subsection{5. Reassess Variance and Mean Consistency, and
ACF}\label{reassess-variance-and-mean-consistency-and-acf}

The series do not have obvious trend or seasonality on the values, and
pattern on the variation.

\begin{Shaded}
\begin{Highlighting}[]
\FunctionTok{par}\NormalTok{(}\AttributeTok{mfrow=}\FunctionTok{c}\NormalTok{(}\DecValTok{1}\NormalTok{,}\DecValTok{1}\NormalTok{))}
\FunctionTok{plot}\NormalTok{(working\_series, }\AttributeTok{main=}\NormalTok{working\_series\_name,}\AttributeTok{type=}\StringTok{"l"}\NormalTok{, }\AttributeTok{ylab=}\NormalTok{working\_ylabel) }
\FunctionTok{abline}\NormalTok{(}\AttributeTok{h=}\DecValTok{0}\NormalTok{, }\AttributeTok{col=}\DecValTok{2}\NormalTok{, }\AttributeTok{lwd=}\DecValTok{2}\NormalTok{)}
\FunctionTok{abline}\NormalTok{(}\AttributeTok{h=}\FunctionTok{mean}\NormalTok{(working\_series), }\AttributeTok{col=}\DecValTok{3}\NormalTok{, }\AttributeTok{lty=}\DecValTok{2}\NormalTok{)}
\end{Highlighting}
\end{Shaded}

\pandocbounded{\includegraphics[keepaspectratio]{CamposCastilloChico_HW1_files/figure-latex/unnamed-chunk-31-1.pdf}}

The variance and mean still seem to differ across years.

\begin{Shaded}
\begin{Highlighting}[]
\FunctionTok{boxplot}\NormalTok{(working\_series }\SpecialCharTok{\textasciitilde{}} \FunctionTok{floor}\NormalTok{(}\FunctionTok{time}\NormalTok{(working\_series)), }
        \AttributeTok{main=}\FunctionTok{paste}\NormalTok{(}\StringTok{"Boxplots by Year:"}\NormalTok{, working\_series\_name),}
        \AttributeTok{xlab=}\StringTok{"Year"}\NormalTok{, }\AttributeTok{ylab=}\NormalTok{working\_ylabel)}
\end{Highlighting}
\end{Shaded}

\pandocbounded{\includegraphics[keepaspectratio]{CamposCastilloChico_HW1_files/figure-latex/unnamed-chunk-32-1.pdf}}

ACF has fast decay which suggests stationarity. PACF shows high spikes
at lag 1 and lag 2. We can still try to improve the stationarity by
performing 1-lag differencing.

\begin{Shaded}
\begin{Highlighting}[]
\CommentTok{\# ACF/PACF to check for non{-}stationarity}
\FunctionTok{par}\NormalTok{(}\AttributeTok{mfrow=}\FunctionTok{c}\NormalTok{(}\DecValTok{1}\NormalTok{,}\DecValTok{2}\NormalTok{))}
\FunctionTok{acf}\NormalTok{(working\_series, }\AttributeTok{ylim=}\FunctionTok{c}\NormalTok{(}\SpecialCharTok{{-}}\DecValTok{1}\NormalTok{,}\DecValTok{1}\NormalTok{), }\AttributeTok{lag.max=}\FunctionTok{max}\NormalTok{(}\DecValTok{60}\NormalTok{, }\DecValTok{3}\SpecialCharTok{*}\NormalTok{freq), }
    \AttributeTok{lwd=}\DecValTok{2}\NormalTok{, }\AttributeTok{main=}\FunctionTok{paste}\NormalTok{(}\StringTok{"ACF:"}\NormalTok{, working\_series\_name),}
\NormalTok{)}
\FunctionTok{pacf}\NormalTok{(working\_series, }\AttributeTok{ylim=}\FunctionTok{c}\NormalTok{(}\SpecialCharTok{{-}}\DecValTok{1}\NormalTok{,}\DecValTok{1}\NormalTok{), }\AttributeTok{lag.max=}\FunctionTok{max}\NormalTok{(}\DecValTok{60}\NormalTok{, }\DecValTok{3}\SpecialCharTok{*}\NormalTok{freq), }
     \AttributeTok{lwd=}\DecValTok{2}\NormalTok{, }\AttributeTok{main=}\FunctionTok{paste}\NormalTok{(}\StringTok{"PACF:"}\NormalTok{, working\_series\_name),}
     \AttributeTok{col=}\FunctionTok{c}\NormalTok{(}\DecValTok{2}\NormalTok{, }\FunctionTok{rep}\NormalTok{(}\DecValTok{1}\NormalTok{, }\DecValTok{59}\NormalTok{)))}
\end{Highlighting}
\end{Shaded}

\pandocbounded{\includegraphics[keepaspectratio]{CamposCastilloChico_HW1_files/figure-latex/acf_pacf_current-1.pdf}}

\begin{Shaded}
\begin{Highlighting}[]
\FunctionTok{par}\NormalTok{(}\AttributeTok{mfrow=}\FunctionTok{c}\NormalTok{(}\DecValTok{1}\NormalTok{,}\DecValTok{1}\NormalTok{))}
\end{Highlighting}
\end{Shaded}

First order differencing does not increase the variance so we can still
proceed with it.

\begin{Shaded}
\begin{Highlighting}[]
\FunctionTok{formatC}\NormalTok{(}\FunctionTok{var}\NormalTok{(working\_series), }\AttributeTok{format =} \StringTok{"f"}\NormalTok{, }\AttributeTok{digits =} \DecValTok{7}\NormalTok{)}
\end{Highlighting}
\end{Shaded}

\begin{verbatim}
## [1] "0.0073107"
\end{verbatim}

\begin{Shaded}
\begin{Highlighting}[]
\FunctionTok{formatC}\NormalTok{(}\FunctionTok{var}\NormalTok{(}\FunctionTok{diff}\NormalTok{(working\_series)), }\AttributeTok{format =} \StringTok{"f"}\NormalTok{, }\AttributeTok{digits =} \DecValTok{7}\NormalTok{)}
\end{Highlighting}
\end{Shaded}

\begin{verbatim}
## [1] "0.0042113"
\end{verbatim}

\subsection{6. Perform 1 lag
differencing}\label{perform-1-lag-differencing}

We apply 1-d differencing.

\begin{Shaded}
\begin{Highlighting}[]
\NormalTok{lag }\OtherTok{=} \DecValTok{1}
\NormalTok{acf\_values }\OtherTok{\textless{}{-}} \FunctionTok{acf}\NormalTok{(working\_series, }\AttributeTok{plot=}\ConstantTok{FALSE}\NormalTok{, }\AttributeTok{lag.max=}\NormalTok{freq)}\SpecialCharTok{$}\NormalTok{acf}
\NormalTok{d1d12series }\OtherTok{\textless{}{-}} \FunctionTok{diff}\NormalTok{(working\_series, }\AttributeTok{lag=}\NormalTok{lag)}
\end{Highlighting}
\end{Shaded}

\begin{Shaded}
\begin{Highlighting}[]
\NormalTok{working\_series\_name }\OtherTok{\textless{}{-}} \FunctionTok{paste0}\NormalTok{(}\StringTok{"Δ"}\NormalTok{, lag,}\StringTok{" "}\NormalTok{,working\_series\_name)}
\NormalTok{working\_ylabel }\OtherTok{\textless{}{-}} \FunctionTok{paste0}\NormalTok{(}\StringTok{"Δ"}\NormalTok{, lag, }\StringTok{" "}\NormalTok{,working\_ylabel)}
\NormalTok{working\_series }\OtherTok{\textless{}{-}}\NormalTok{ d1d12series}
\end{Highlighting}
\end{Shaded}

\subsubsection{6.1 Reassess Variance and Mean Consistency, and
ACF}\label{reassess-variance-and-mean-consistency-and-acf-1}

\begin{Shaded}
\begin{Highlighting}[]
\FunctionTok{par}\NormalTok{(}\AttributeTok{mfrow=}\FunctionTok{c}\NormalTok{(}\DecValTok{1}\NormalTok{,}\DecValTok{1}\NormalTok{))}
\FunctionTok{plot}\NormalTok{(working\_series, }\AttributeTok{main=}\NormalTok{working\_series\_name,}\AttributeTok{type=}\StringTok{"l"}\NormalTok{, }\AttributeTok{ylab=}\NormalTok{working\_ylabel) }
\FunctionTok{abline}\NormalTok{(}\AttributeTok{h=}\DecValTok{0}\NormalTok{, }\AttributeTok{col=}\DecValTok{2}\NormalTok{, }\AttributeTok{lwd=}\DecValTok{2}\NormalTok{)}
\FunctionTok{abline}\NormalTok{(}\AttributeTok{h=}\FunctionTok{mean}\NormalTok{(working\_series), }\AttributeTok{col=}\DecValTok{3}\NormalTok{, }\AttributeTok{lty=}\DecValTok{2}\NormalTok{)}
\end{Highlighting}
\end{Shaded}

\pandocbounded{\includegraphics[keepaspectratio]{CamposCastilloChico_HW1_files/figure-latex/unnamed-chunk-36-1.pdf}}

The mean seem to be constant and the variance closer with each other.

\begin{Shaded}
\begin{Highlighting}[]
\FunctionTok{boxplot}\NormalTok{(working\_series }\SpecialCharTok{\textasciitilde{}} \FunctionTok{floor}\NormalTok{(}\FunctionTok{time}\NormalTok{(working\_series)), }
        \AttributeTok{main=}\FunctionTok{paste}\NormalTok{(}\StringTok{"Boxplots by Year:"}\NormalTok{, working\_series\_name),}
        \AttributeTok{xlab=}\StringTok{"Year"}\NormalTok{, }\AttributeTok{ylab=}\NormalTok{working\_ylabel)}
\end{Highlighting}
\end{Shaded}

\pandocbounded{\includegraphics[keepaspectratio]{CamposCastilloChico_HW1_files/figure-latex/unnamed-chunk-37-1.pdf}}

Both ACF and PACF suggests stationarity because of the rapid decay and
less number of strong PACF. This can be modeled using ARMA.

\begin{Shaded}
\begin{Highlighting}[]
\CommentTok{\# ACF/PACF to check for non{-}stationarity}
\FunctionTok{par}\NormalTok{(}\AttributeTok{mfrow=}\FunctionTok{c}\NormalTok{(}\DecValTok{1}\NormalTok{,}\DecValTok{2}\NormalTok{))}
\FunctionTok{acf}\NormalTok{(working\_series, }\AttributeTok{ylim=}\FunctionTok{c}\NormalTok{(}\SpecialCharTok{{-}}\DecValTok{1}\NormalTok{,}\DecValTok{1}\NormalTok{), }\AttributeTok{lag.max=}\FunctionTok{max}\NormalTok{(}\DecValTok{60}\NormalTok{, }\DecValTok{3}\SpecialCharTok{*}\NormalTok{freq), }
    \AttributeTok{lwd=}\DecValTok{2}\NormalTok{, }\AttributeTok{main=}\FunctionTok{paste}\NormalTok{(}\StringTok{"ACF:"}\NormalTok{, working\_series\_name),}
\NormalTok{)}
\FunctionTok{pacf}\NormalTok{(working\_series, }\AttributeTok{ylim=}\FunctionTok{c}\NormalTok{(}\SpecialCharTok{{-}}\DecValTok{1}\NormalTok{,}\DecValTok{1}\NormalTok{), }\AttributeTok{lag.max=}\FunctionTok{max}\NormalTok{(}\DecValTok{60}\NormalTok{, }\DecValTok{3}\SpecialCharTok{*}\NormalTok{freq), }
     \AttributeTok{lwd=}\DecValTok{2}\NormalTok{, }\AttributeTok{main=}\FunctionTok{paste}\NormalTok{(}\StringTok{"PACF:"}\NormalTok{, working\_series\_name),}
     \AttributeTok{col=}\FunctionTok{c}\NormalTok{(}\DecValTok{2}\NormalTok{, }\FunctionTok{rep}\NormalTok{(}\DecValTok{1}\NormalTok{, }\DecValTok{59}\NormalTok{)))}
\end{Highlighting}
\end{Shaded}

\pandocbounded{\includegraphics[keepaspectratio]{CamposCastilloChico_HW1_files/figure-latex/acf_pacf_final-1.pdf}}

\begin{Shaded}
\begin{Highlighting}[]
\FunctionTok{par}\NormalTok{(}\AttributeTok{mfrow=}\FunctionTok{c}\NormalTok{(}\DecValTok{1}\NormalTok{,}\DecValTok{1}\NormalTok{))}
\end{Highlighting}
\end{Shaded}

The variance increased if we perform one more differencing so there is
really no need to do one more.

\begin{Shaded}
\begin{Highlighting}[]
\FunctionTok{formatC}\NormalTok{(}\FunctionTok{var}\NormalTok{(working\_series), }\AttributeTok{format =} \StringTok{"f"}\NormalTok{, }\AttributeTok{digits =} \DecValTok{7}\NormalTok{)}
\end{Highlighting}
\end{Shaded}

\begin{verbatim}
## [1] "0.0042113"
\end{verbatim}

\begin{Shaded}
\begin{Highlighting}[]
\FunctionTok{formatC}\NormalTok{(}\FunctionTok{var}\NormalTok{(}\FunctionTok{diff}\NormalTok{(working\_series)), }\AttributeTok{format =} \StringTok{"f"}\NormalTok{, }\AttributeTok{digits =} \DecValTok{7}\NormalTok{)}
\end{Highlighting}
\end{Shaded}

\begin{verbatim}
## [1] "0.0123498"
\end{verbatim}

Overall, after applying a logarithmic transformation, seasonal
differencing, and 1-period differencing, the AIRBCN series appears to be
weakly stationary with stabilized variance and mean.

\end{document}
